% concatenated from Final_LZ.tex
\documentclass[11pt,reqno]{amsart}
\usepackage[T1]{fontenc}
\usepackage{lmodern,amsmath,amssymb,amsthm,mathtools,mathrsfs,esint}
\usepackage{enumitem,booktabs,array,needspace}
\usepackage[margin=1in]{geometry}
\usepackage[expansion=false]{microtype}
\usepackage[colorlinks=true,linkcolor=blue,citecolor=blue,urlcolor=blue]{hyperref}
\hypersetup{pdftitle={Non-simple blow-up for the Chern--Simons--Higgs equation: A priori analysis and constructions},
pdfauthor={Youngae Lee, and Lei Zhang},
pdfsubject={A priori estimates and constructions in both profile regimes}}
\allowdisplaybreaks[2]
\AtBeginDocument{%
  \setlength{\abovedisplayskip}{3.5pt plus 3pt minus 1pt}%
  \setlength{\belowdisplayskip}{3.5pt plus 3pt minus 1pt}%
  \setlength{\abovedisplayshortskip}{0pt plus 2pt}%
  \setlength{\belowdisplayshortskip}{2pt plus 2pt minus 1pt}%
  \setlength{\jot}{2pt}%
}
\numberwithin{equation}{section}
\newtheorem{theorem}{Theorem}[section]
\newtheorem{proposition}[theorem]{Proposition}
\newtheorem{lemma}[theorem]{Lemma}
\newtheorem{corollary}[theorem]{Corollary}
\theoremstyle{definition}
\newtheorem{definition}[theorem]{Definition}
\newtheorem{assumption}[theorem]{Assumption}
\theoremstyle{remark}
\newtheorem{remark}[theorem]{Remark}
\newtheorem*{remark*}{Remark}
\newcommand{\R}{\mathbb R}
\newcommand{\T}{\mathbb T}
\newcommand{\cS}{\mathcal S}
\newcommand{\eps}{\varepsilon}
\newcommand{\etae}{\eta_{\eps}}
\newcommand{\de}{\delta_{\eps}}
\newcommand{\we}{w_{\eps}}
\newcommand{\tr}{\operatorname{tr}}
\newcommand{\dist}{\operatorname{dist}}
\newcommand{\loc}{\mathrm{loc}}
\newcommand{\comp}{\mathrm{comp}}
\newcommand{\tot}{\mathrm{tot}}

\newcommand{\C}{\mathbb C}
\newcommand{\dd}{\,dx}
\newcommand{\calA}{\mathcal A}
\DeclareMathOperator{\Res}{Res}
\DeclareMathOperator{\osc}{osc}
\title[Non-simple blow-up for the Chern--Simons--Higgs equation]
{Non-simple blow-up for the Chern--Simons--Higgs equation:
A priori analysis and constructions}
 
\author{Youngae Lee}
\address[Youngae Lee]{Department of Mathematical Sciences,
Ulsan National Institute of Science and Technology (UNIST), Ulsan, Republic of Korea}
\email{youngaelee@unist.ac.kr}
\author{Lei Zhang}
\address[Lei Zhang]{Department of Mathematics, University of Florida,
Gainesville, FL 32611, USA}
\email{leizhang@ufl.edu} 
\date{}
\subjclass[2020]{35B44, 35J61, 35J75}
\keywords{Chern--Simons--Higgs equation, non-simple blow-up, a priori estimate,
finite-height profile, mean-field limit, collapsing bubbles, Pohozaev identity}
\begin{document}
 
\begin{abstract}

 Blow-up in the self-dual Chern--Simons--Higgs equation has a rich
multiscale structure, but the detailed theory has largely focused on
the simple regime, where one entire profile resolves each
concentration point. By contrast, non-simple blow-up is analytically more delicate:
several interacting regular cores must be resolved simultaneously
as they collapse toward a single vortex while remaining separated
on their own scales. Because of this difficulty, the existence of
such solutions has remained a long-standing open problem. To the best of our
knowledge, we prove the first existence result for non-simple blow-up
solutions of this equation, covering in particular the previously
unresolved finite-height Chern--Simons regime.

We also show that non-simple clusters have equal
limiting core masses, a common profile type, and regular-polygon
arrangements. We determine their possible total masses and derive a
mass-gap criterion that rules out such blow-up for certain vortex
configurations. Our quantitative profile and moment estimates for general
finite-height clusters guide the design of the approximate solution
and clarify the core interactions underlying the finite-height
construction. A separate
construction on the unit disk produces two mean-field cores along a
sequence with varying Dirichlet traces.
\end{abstract}

\maketitle

\section{Introduction and main results}\label{sec:introduction}

On a smooth closed connected Riemannian surface $(M,g)$, the scalar
self-dual Chern--Simons--Higgs equation is
\begin{equation}\label{eq:main}
 \Delta_gu_\eps+\eps^{-2}F(u_\eps)
 =4\pi\sum_{i=1}^{N_{\mathrm v}}N_i\delta_{p_i},
\end{equation}
where $\eps>0$, the vortex points $p_i$ and strengths $N_i>0$ are fixed,
$N_{\tot}=\sum_iN_i$, and
\begin{equation}\label{eq:F-definition}
 F(t)=e^t(1-e^t),\qquad \mathcal P(t)=e^t-\tfrac12e^{2t},\qquad \mathcal P'=F.
\end{equation}
For integer $N_i$, the function $u_\eps$ is the logarithm of the squared
Higgs modulus, normalized to one in the vacuum; its singularity
$2N_i\log d_g(x,p_i)$ gives the source $4\pi N_i\delta_{p_i}$.
The equation describes electrically charged vortices, with a flat
torus representing a periodic condensate cell
\cite{HongKimPac,JackiwWeinberg,SpruckYang}.  
For fixed nonempty   vortex data on a flat torus, solutions
exist precisely for $0<\eps\le\eps_c$, with at least two solutions for
$0<\eps<\eps_c$ \cite{CaffarelliYang,Tarantello}.


Since the seminal vortex constructions, the self-dual Chern--Simons--Higgs equation has provided a fertile setting for nonlinear elliptic analysis: existence and multiplicity, concentration, local profiles, and quantized mass have all been studied extensively; representative contributions include \cite{CaffarelliYang,Tarantello,CFL,ChoeKim,CKL,LinYanCMP,LinYanARMA,DelPinoEspositoFigueroaMusso,LinYanAIHP,LinYanAdvMath}. In the blow-up theory, however, the most complete results concern simple concentration, for which a single entire profile captures the local behavior. A non-simple cluster requires several regular cores to be resolved simultaneously while their mutual distances vanish on the ambient scale. This creates interacting translation, scale, and mass modes, and prevents a direct reduction to the established one-bubble analysis.  {Non-simple Liouville-type blow-up has been analyzed, and in several settings constructed, for singular Liouville and mean-field equations \cite{KuoLin,BartolucciTarantello,DAprileWeiZhangDCDS,DAprileWeiZhangCVPDE,DAprileWeiZhangAJM}. The corresponding existence problem for the full Chern--Simons--Higgs equation has nevertheless remained open in both profile regimes.} To the best of our knowledge, the present paper gives the first existence result for non-simple blow-up solutions of this equation; in particular, it gives the first construction in the finite-height Chern--Simons regime.  



 

\subsection{Concentration regimes and related work}

A concentrating regular core has a profile solving either
\[
 \Delta U+e^U(1-e^U)=0
 \quad\hbox{or}\quad \Delta V+e^V=0
 \qquad\hbox{in }\R^2.
\]
The finite-height Chern--Simons profile has height-dependent mass
greater than $8\pi$; the mean-field Liouville profile has mass $8\pi$
\cite{CFL,Choe,ChoeKim,CKL,LinYanAIHP}. Thus the finite-height
interaction problem must determine the individual masses as well
as the core positions.

Lin--Yan~\cite{LinYanAIHP} analyze both regimes, exclude their
coexistence when all blow-up points avoid the vortices, and construct
multiple Chern--Simons bubbles. Their periodic constructions also
include strength-one vortex bubbling~\cite{LinYanCMP} and
mean-field bubbles~\cite{LinYanARMA}; necessary conditions and local
uniqueness appear in~\cite{LinYanAdvMath}.
Del Pino--Esposito--Figueroa--Musso~\cite{DelPinoEspositoFigueroaMusso}
construct periodic mean-field condensates using weighted entire
Liouville profiles. Although the cited constructions allow multiple bubbles, each
concentration point in the constructed solutions is described by a
single local profile. However, the existence of non-simple solutions,
in which several regular cores collapse at one vortex while remaining
separated on their own scales, had remained open.

Li--Liu~\cite[Theorems~1.2 and~1.3]{LiLiu} prove an energy identity
without neck loss at vortex points, describe the singular and regular
bubble alternatives, and obtain simplicity in low-vorticity cases. In
this paper, we complement this analysis by determining the individual
core masses and relative positions, and by establishing the
realizability of non-simple clusters.

For singular Liouville and mean-field equations, the regular-polygon
structure and pointwise estimates were obtained independently by
Kuo--Lin~\cite{KuoLin} and Bartolucci--Tarantello~\cite[Theorem~1.4
and Section~6]{BartolucciTarantello}. Higher-order estimates and
vanishing results appear in
\cite{WeiZhang,WeiZhangVanishing,WeiZhangLaplacian,DAprileWeiZhangTAMS},
and construction or exclusion results in
\cite{DAprileWeiZhangDCDS,DAprileWeiZhangCVPDE,DAprileWeiZhangAJM}. 

The new difficulty in the finite-height regime is to control
interacting cores as they collapse at a single vortex, while keeping
the structural and local estimates uniform across scales.  Our quantitative result, Theorem~\ref{prop:refined-profile}, combines
explicit center-shift asymptotics, a second Fourier correction, and
a controlled comparison of local and radial-profile masses for
individual finite-height cores. Its proof and moment formulas are
in Appendix~\ref{app:pointwise-proofs}. The structural results cover
both regimes; a complementary disk construction realizes the
mean-field regime with varying boundary traces.

\subsection{Cluster structure and a mass-gap criterion}

We first specify the meaning of non-simple blow-up. Let
$\mathcal Z=\{p_1,\ldots,p_{N_{\mathrm v}}\}$, and write $|M|_g$
for the area of $M$. For a concentrating family $(u_\eps)$ of solutions
of \eqref{eq:main} in the setting of Section~\ref{sec:setting},
we use the following normalization.
\begin{align}
 -\Delta_{g,x}G_g(x,y)&=\delta_y-|M|_g^{-1},\qquad
 \int_MG_g(x,y)\,dV_g(x)=0,\label{eq:green}\\
 u_0(x)&=-4\pi\sum_iN_iG_g(x,p_i),\qquad
 w_\eps=u_\eps-u_0-2\log\eps.\label{eq:u0-w}
\end{align}
Here $d_g$ denotes geodesic distance, and $B_r^g(q)$ denotes
its corresponding geodesic disk. In local conformal coordinates
and in rescaled variables, $B_r(q)$ denotes a Euclidean disk.
\begin{definition}[Blow-up and non-simplicity]\label{def:simple}
A point $q$ belongs to the blow-up set $\cS_u$ if
$w_\eps(x_{\eps,q})\to+\infty$ for some $x_{\eps,q}\to q$.
Put $n_q=0$ when $q\notin\mathcal Z$, and $n_q=N_i$ when $q=p_i$.
For a vortex, set $\xi_{\eps,q}=q$; for a regular blow-up point,
choose a local maximum $\xi_{\eps,q}\to q$ of $w_\eps$.
The point is simple if such a choice satisfies, for some small $r>0$,
\begin{equation}\label{eq:simple}
 \limsup_{\eps\to0}\sup_{B_r(q)}
 \{w_\eps(x)+2(1+n_q)\log d_g(x,\xi_{\eps,q})\}<\infty;
\end{equation}
otherwise it is non-simple, after passing to a subsequence.
\end{definition}
The moving center at a regular point allows a single core to drift
toward its limit; at a vortex the singular source fixes the center.

Normalize masses by $2\pi$, so the total mass is $2N_{\tot}$.
Under the hypotheses of Section~\ref{sec:setting}, logarithmic
comparisons and exterior estimates based on \cite{CKL,LinYanAIHP}
give a common limiting core mass $m\ge4$: all cores are mean-field
if $m=4$ and Chern--Simons if $m>4$, without prescribing their type
or a common height. In Section~\ref{sec:cluster}, Pohozaev identities
and a further rescaling yield a regular $S$-gon and
\[
 m(S-1)=4N_i,\qquad M_{p_i}=mS=\frac{4N_iS}{S-1}.
\]
In particular, $1\le S-1\le N_i$, with equality on the right precisely in the
mean-field case. The finite-height hypothesis used for the sharp
estimates in Section~\ref{sec:pointwise} is not imposed in this
structural result.

\begin{theorem}[Non-simple masses and the four-unit gap criterion]
\label{thm:ratio-free-four-unit-gap}
Let \((u_\eps)\) be a concentrating non-topological family of solutions of
\eqref{eq:main} in the setting of Section~\ref{sec:setting}.
For a singular source $p_i$, put
\[
 \mathscr A_i:=\{k\in\mathbb N:1\le k\le N_i\},
 \qquad
 M_{i,k}:=\frac{4N_i(k+1)}{k}.
\]
If a non-simple cluster occurs at $p_i$ with $S=k+1$ cores, then
$k\in\mathscr A_i$, its normalized mass is $M_{i,k}$, and necessarily
\begin{equation}\label{eq:necessary-alternative}
2N_{\mathrm{tot}}-M_{i,k}\in\{0\}\cup[4,\infty).
\end{equation}
In particular, suppose that, for every $i$ and every $k\in\mathscr A_i$,
\begin{equation}
 2N_{\rm tot}-M_{i,k}\in(-\infty,0)\cup(0,4).
 \label{eq:ratio-free-four-unit-gap}
\end{equation}
Then every blow-up point is simple.  
\end{theorem}

An additional blow-up point has mass at least four, which excludes
a positive remainder below four in either regime without a further
scale-rate assumption. The zero-remainder alternative is attained
below: two cores require $N_{\tot}=4N_i$, and the construction has
$N_i=2$, $N_{\tot}=8$.

\subsection{Existence of a non-simple Chern--Simons cluster}

Our first existence result realizes the resonant finite-height configuration
on the square flat torus. Let
\(
 \T=\R^2/\mathbb Z^2,\) \(a=(1/2,0),\) \( e=(0,1).
\)
Let \(G\) be its Green function with normalization
\begin{equation}\label{cs:eq:green}
 -\Delta_xG(x,p)=\delta_p-1,\qquad \int_\T G(x,p)\,dx=0,
\end{equation}
and use the specialization of the background in \eqref{eq:u0-w},
\begin{equation}\label{cs:eq:u0}
 u_0=-8\pi G(\cdot,0)-24\pi G(\cdot,a).
\end{equation}
The vortex configuration \(2[0]+6[a]\) is fixed throughout the construction.
Let \(U\) be the radial entire profile centered at the origin with
\[
 \Delta U+F(U)=0\quad\text{in }\R^2,\qquad U<0,\qquad
 \frac1{2\pi}\int_{\R^2}F(U)\,dx=8.
\]
\begin{theorem}[Non-simple blow-up of Chern--Simons type on the square torus]\label{cs:thm:existence}
There is a constant $d_*>0$ such that, for every sufficiently small $\eps>0$,
there is a distributional solution
\begin{equation}\label{cs:eq:target}
 \Delta u_\eps+\eps^{-2}F(u_\eps)
       =8\pi\delta_0+24\pi\delta_a\quad\hbox{on }\T.
\end{equation}
It satisfies $u_\eps-u_0\in C^\infty(\T)$ and is invariant under
$(x,y)\mapsto(-x,y)$ and $(x,y)\mapsto(x,-y)$.
There are $\gamma>0$ and local maxima
$q_{\eps,\pm}=\pm\delta_\eps e$ such that
\begin{align}
 &\sup_\T u_\eps\le-\gamma,\qquad
 \delta_\eps=d_*\eps^{2/3}(1+o(1)),\label{cs:eq:scale}\\
 &u_\eps(q_{\eps,\pm}+\eps z)\longrightarrow U(z)
             \quad\hbox{in }C^2_{\rm loc}(\R^2),\label{cs:eq:profiles}\\
 &\eps^{-2}F(u_\eps)\dd\rightharpoonup32\pi\delta_0.\label{cs:eq:measure}
\end{align}
Each limiting core has mass $16\pi$, and there is no missing mass.
For $w_\eps=u_\eps-u_0-2\log\eps$,
\[
 w_\eps\longrightarrow-\infty\quad\hbox{locally uniformly on }\T\setminus\{0\},
 \qquad
 w_\eps(q_{\eps,+})+6\log|q_{\eps,+}|\longrightarrow+\infty .
\]
In particular, this is non-simple blow-up of finite-height
Chern--Simons type at the strength-two vortex.
\end{theorem}
The two core widths are of order $\eps$, whereas their distance is
asymptotic to $2d_*\eps^{2/3}$, so the cores remain distinct on their
own scales. Their combined mass exhausts the total mass and realizes
the zero-remainder alternative in \eqref{eq:necessary-alternative}.
Among configurations with positive integer multiplicities $N_i$, the
example has the smallest total vorticity compatible with a finite-height
non-simple cluster: $N_{\tot}=8$. Indeed, $1\le k<N_i$ then implies
$N_i\ge k+1$, hence $M_{i,k}\ge4(k+1)^2/k\ge16$ and $2N_{\tot}\ge16$.
The two-core configuration is also compatible with the quadratic
mass identity in \cite[Theorem~1.2]{LiLiu}: for normalized core masses
$b_1=b_2=8$ at a strength-two vortex, that identity reads
\[
 b_1^2+b_2^2=(b_1+b_2)^2-4N_i(b_1+b_2),
 \qquad 8^2+8^2=16^2-8\cdot16.
\]
Thus the construction realizes an equality permitted by that identity.

To the best of our knowledge, Theorem~\ref{cs:thm:existence} provides
the first existence result for non-simple blow-up with finite-height
Chern--Simons profiles. The separation law distinguishes it from simple blow-up
described by one entire profile. Here the core width $\eps$ and the
half-separation $\delta$ satisfy $\eps/\delta\to0$: no single rescaling
both resolves the cores and keeps their centers at distinct finite
positions. Estimates must remain uniform as the centers merge, a
degeneration not addressed directly by constructions with distinct
limiting centers.

Moreover, the leading vortex--core and inter-core forces cancel.
Isotropic second moments vanish against the harmonic interaction field;
anisotropic moments supply the next stress contribution. With
$\delta=d\eps^{2/3}$, the reduced equation is
\[
 \lambda d-\frac{\mathfrak b}{2d^5}+o(1)=0,
 \qquad \lambda>0,\quad \mathfrak b>0.
\]
The torus Green function determines $\lambda$, and the linearized
second angular mode determines $\mathfrak b$; see
\eqref{cs:eq:lambda} and \eqref{cs:eq:beta}. Uniform control of the
collapsing radial and logarithmic modes, followed by the continuous
translation reduction, gives a solution for every small $\eps$ and
$d_*=(\mathfrak b/(2\lambda))^{1/6}$.
\subsection{Existence of a non-simple mean-field cluster}

To complement the finite-height construction, we realize the mean-field
endpoint with two cores. This second result concerns the same scalar
nonlinearity on the
Euclidean unit disk \(B_1\subset\R^2\), with one fixed strength-one vortex:
\begin{equation}\label{mf:eq:CS}
\Delta u+\eps^{-2}F(u)=4\pi\delta_0
\qquad\text{in }B_1\subset\R^2.
\end{equation}
We identify \(\R^2\) with \(\C\) by \(z=x+iy\). For a sequence with couplings
\(\eps_k\to0\), write
\(w_k=u_k-2\log\eps_k-2\log|z|.\)
In this setting non-simplicity at the origin means that, for every
sufficiently small fixed $r>0$,
\begin{equation}\label{mf:eq:nonsimplicity}
 \limsup_{k\to\infty}\sup_{0<|z|<r}
       \{w_k(z)+4\log|z|\}=+\infty.
\end{equation}
The Dirichlet traces are part of the construction and are specified
in the theorem.
\begin{theorem}[Non-simple blow-up of mean-field type on the unit disk]\label{mf:thm:main}
There exist $s_k\downarrow0$, $\eps_k\downarrow0$, and solutions $u_k$
of \eqref{mf:eq:CS}, with $u_k-2\log|z|\in C^\infty(\overline{B_1})$,
having the following properties:
\begin{enumerate}
\item The prescribed boundary values are
\begin{equation}\label{mf:eq:boundary}
u_k(e^{i\theta})=
\log\frac{32\eps_k^2s_k^4}
{\bigl(s_k^4+1+s_k^2-2s_k\cos(2\theta)\bigr)^2}.
\end{equation}
Their oscillation is $O(s_k)$.
\item $\sup_{B_1}u_k\to-\infty$.
\item With
\(
D_k=\eps_k^{-2}F(u_k),
\)
we have $D_k\dd\rightharpoonup16\pi\delta_0$ as finite measures
on $\overline{B_1}$.
\item Put $q_{\pm,k}=\pm\sqrt{s_k}$ and $\ell_k=s_k^{3/2}/2$.
For each sign there is a local maximum of $u_k$ at
$q_{\pm,k}+o(\ell_k)$, and
$
\lim_{R\to\infty}\lim_{k\to\infty}
\int_{B_{R\ell_k}(q_{\pm,k})}D_k\dd=8\pi.
$
The origin is the unique blow-up point of $w_k$ and is non-simple
in the sense of \eqref{mf:eq:nonsimplicity}.
\end{enumerate}
\end{theorem}

\begin{remark}\label{mf:rem:scope}
The traces in \eqref{mf:eq:boundary} vary with $k$.
The coupling sequence is chosen diagonally; the theorem asserts neither
a universal power law between $\eps_k$ and $s_k$ nor existence for every
sufficiently small coupling with these traces.
\end{remark}

Sections~\ref{sec:setting}--\ref{sec:cluster} develop the cluster
structure; Section~\ref{sec:pointwise} states the quantitative estimates,
proved with their moment formulas in Appendix~\ref{app:pointwise-proofs}.
Sections~\ref{sec:cs-construction} and~\ref{sec:mf-construction}
give the periodic and disk constructions.

\section{Geometric setting and profile regimes}\label{sec:setting}
We use the notation and definition of blow-up introduced in
Section~\ref{sec:introduction}. For the a priori analysis, consider
a sequence of solutions of \eqref{eq:main} satisfying
\( \limsup_{\eps\to0}\sup_Mu_\eps<0.
\)
We assume the concentrating alternative of \cite{ChoeKim,LinYanAIHP}:
$\cS_u$ is finite and nonempty, and, after extraction,
\begin{gather}
 \mu_\eps:=\frac{F(u_\eps)}{2\pi\eps^2}\,dV_g
 \rightharpoonup\sum_{q\in\cS_u}M_q\delta_q,\qquad M_q\ge4,
 \label{eq:measure-limit}\\
 w_\eps\longrightarrow-\infty
 \quad\text{locally uniformly on }M\setminus\cS_u.\notag
\end{gather}
Integration of \eqref{eq:main} gives
\begin{equation}\label{eq:global-mass}
 \mu_\eps(M)=\sum_{q\in\cS_u}M_q=2N_{\tot}.
\end{equation}
With $\overline u_\eps=|M|_g^{-1}\int_Mu_\eps\,dV_g$, the exact
Green representation in this normalization is
\begin{equation}\label{eq:global-green-representation}
 u_\eps(x)-u_0(x)=\overline u_\eps
                +2\pi\int_M G_g(x,y)\,d\mu_\eps(y).
\end{equation}
Indeed, $-\Delta_g(u_\eps-u_0)=\eps^{-2}F(u_\eps)
-4\pi N_{\tot}/|M|_g$, and $u_0$ has zero average.
For an iterated error $o_R(1)+o_\eps(1)$, the limit in $\eps$ is
taken first at fixed $R$, followed by $R\to\infty$, unless an
explicit uniform bound is stated.

All masses $m$, $\widetilde m$, and $M_q$ are normalized by $2\pi$;
$M$ denotes the surface. A tilde marks actual finite-disk masses,
while $m(\alpha)$ is the profile mass at height $\alpha$.
The small exponent $\sigma$ is used in H\"older norms and sublinear weights.
Regular blow-up points are simple \cite{ChoeKim,LinYanAIHP}.
The finite-height hypothesis is used in Section~\ref{sec:pointwise}
and Appendix~\ref{app:pointwise-proofs}; Section~\ref{sec:cluster}
covers both regimes.

\subsection{Local coefficients}
At $p_i$, choose a conformal normal coordinate with
\begin{equation}\label{eq:isothermal}
 g=e^{\varphi_i(x)}|dx|^2,\quad \varphi_i(0)=0,\quad
 \nabla\varphi_i(0)=0,\quad D^2\varphi_i(0)=-K_g(p_i)I_2.
\end{equation}
The harmonic quadratic part is removed by a holomorphic coordinate
change. Write $\gamma_i(x)=G_g(x,p_i)+(2\pi)^{-1}\log|x|$.
We fix the normalization of the smooth positive coefficient $h_{i,g}$ by
\begin{equation}\label{eq:explicit-fixed-coefficient}
 \log h_{i,g}=4\pi(N_{\tot}-N_i)\gamma_i
               -4\pi\sum_{j\ne i}N_jG_g(\cdot,p_j).
\end{equation}
For the limiting masses in \eqref{eq:measure-limit}, define the field
of the other blow-up points by
\begin{equation}\label{eq:other-core-harmonic-field}
 \Psi_i(x)=2\pi\sum_{q\in\cS_u\setminus\{p_i\}}M_q
 \bigl[G_g(x,q)-\gamma_i(x)-G_g(p_i,q)+\gamma_i(p_i)\bigr].
\end{equation}
Thus $\Psi_i(p_i)=0$; an empty sum is zero. On a chart containing no
other blow-up point, each bracket is smooth and harmonic, because
$\Delta G_g(\cdot,q)=\Delta\gamma_i=e^{\varphi_i}/|M|_g$.
Set
\begin{equation}\label{eq:fixed-K}
 \mathcal K_i=\log h_{i,g},\qquad
 \mathcal K_i^{\comp}=\mathcal K_i+\Psi_i.
\end{equation}
Both coefficients are harmonic locally. All derivatives of this fixed
field are bounded on a smaller chart; its additive normalization
cancels in the rescaled equation.

\subsection{Finite-height profiles}
Let $U_\alpha$ be the radial non-topological entire solution
\begin{equation}\label{eq:entire}
 \Delta U_\alpha+F(U_\alpha)=0
 \quad\text{in }\R^2,\qquad U_\alpha(0)=\alpha<0,
 \quad U_\alpha'(0)=0,
\end{equation}
and set
\begin{equation}\label{eq:mass-map}
 m(\alpha)=\frac1{2\pi}\int_{\R^2}F(U_\alpha)\,dz.
\end{equation}
The profile family satisfies
\begin{equation}\label{eq:profile-facts}
 m(\alpha)>4,\qquad m'(\alpha)>0\qquad(\alpha<0).
\end{equation}
The family depends $C^1$-smoothly on $\alpha$, and $m$ maps
$(-\infty,0)$ onto $(4,\infty)$;
see \cite{CFL} and \cite[pp.~840--841]{CKL}. In the notation of
\cite[Lemma~2.4]{CKL}, with vortex number zero, their height $s$ is
our $\alpha$ and their total mass is $2\pi m(\alpha)$.
The negative logarithmic coefficient of $\partial_\alpha U_\alpha$
is $-m'(\alpha)$; that lemma gives $m'(\alpha)>0$.
\begin{assumption}[Finite-height cluster]\label{ass:profile}
The cluster considered in Section~\ref{sec:pointwise} is of
Chern--Simons type, with normalized heights converging to
$\beta\in(-\infty,0)$.
\end{assumption}
After choosing $I\Subset(-\infty,0)$ containing these heights, fix
\begin{equation}\label{eq:profile-hyp}
 4<m_*<\inf_{\alpha\in I}m(\alpha),\qquad
 \inf_{\alpha\in I}m'(\alpha)>0.
\end{equation}
The strict margin in $m_*$ also controls the convergent finite-disk
masses. Constants below are uniform for $\alpha\in I$.

\section{Entire profiles and the linearized operator}\label{sec:linearized}
Put
\begin{equation}\label{eq:linearized-operator}
 L_\alpha=\Delta+F'(U_\alpha).
\end{equation}
The radial equation gives
\begin{equation}\label{eq:radial-monotonicity}
 (rU_\alpha')'=-rF(U_\alpha),\qquad U_\alpha'(r)<0\quad(r>0).
\end{equation}
With $m=m(\alpha)$ and $c_\infty=c_\infty(\alpha)$, we have as $r\to\infty$,
\begin{align}
 U_\alpha(r)&=-m\log r+c_\infty+O(r^{2-m}), \ \
 U_\alpha'(r)=-m/r+O(r^{1-m}), \ \
 U_\alpha''(r)=m/r^2+O(r^{-m}).
 \label{eq:Usecond-asymptotic}
\end{align}
At the origin, writing $\kappa_\alpha=F(\alpha)>0$,
\begin{equation}\label{eq:kappa-alpha}
 D^2U_\alpha(0)=-\frac{\kappa_\alpha}{2}I_2,\qquad
 U_\alpha'(r)=-\frac{\kappa_\alpha}{2}r+O(r^3).
\end{equation}

\begin{lemma}[Sublinear kernel and normalization]
\label{thm:linearized-classification}
If $L_\alpha\phi=0$ in $\R^2$ and
$|\phi(x)|\le C(1+|x|)^\tau$ for some $0<\tau<1$, then
\begin{equation}\label{eq:kernel-classification}
 \phi=c_0\partial_\alpha U_\alpha
       +c_1\partial_{x_1}U_\alpha+c_2\partial_{x_2}U_\alpha.
\end{equation}
In particular, $\phi(0)=0$ and $\nabla\phi(0)=0$ imply $\phi=0$.
\end{lemma}
\begin{proof}
Expand in angular modes. Each radial coefficient of order $k$ satisfies
\begin{equation}\label{eq:mode-equation}
 h''+r^{-1}h'-k^2r^{-2}h+F'(U_\alpha)h=0.
\end{equation}
For $k=0$, the regular solution is a multiple of
$\partial_\alpha U_\alpha$, whose value at zero is one.
For $k=1$, write $h=wU_\alpha'$; then
$(r(U_\alpha')^2w')'=0$. Regularity at zero makes the integration
constant zero, so $h$ is a multiple of $U_\alpha'$.
For $k\ge2$, put $\psi=-U_\alpha'>0$ and choose the sign of a nonzero
regular solution so that $h(r)\sim a r^k$, $a>0$, at zero. Then
\begin{equation}\label{eq:higher-mode-comparison}
 \bigl(r\psi^2(h/\psi)'\bigr)'=(k^2-1)r^{-1}\psi h.
\end{equation}
Thus $h/\psi$ is positive and increasing until any first zero of $h$,
which is impossible. For large $r$, $r\psi^2(h/\psi)'\ge c>0$;
as $\psi\sim m/r$, this gives $(h/\psi)'\ge c'r$ and $h\ge c''r$.
Such a mode cannot be sublinear.
The value and gradient normalizations eliminate the three remaining
coefficients by \eqref{eq:kappa-alpha}.
\end{proof}

\begin{lemma}[Weighted Green estimate]\label{lem:weighted-green}
Let $G_R$ be the positive Dirichlet Green function of $-\Delta$ in
$B_R$, $R\ge2$. For $0<\sigma<1$, $\rho>0$, and every $x\in B_R$,
\begin{equation}\label{eq:weighted-green}
 \int_{B_R}|G_R(x,y)-G_R(0,y)|(1+|y|)^{-2-\rho}\,dy
 \le C_{\rho,\sigma}(1+|x|)^\sigma.
\end{equation}
Also,
\begin{equation}\label{eq:weighted-green-gradient}
 \int_{B_R}|\nabla_xG_R(x,y)|(1+|y|)^{-2-\rho}\,dy
 \le C_{\rho,\sigma}(1+|x|)^{\sigma-1}.
\end{equation}
The constants are independent of $R$; the same bounds hold for disks
whose centers are $o(1)$ from the normalization point.
\end{lemma}
\begin{proof}
The disk formula
\(
 G_R(x,y)=\frac1{2\pi}\log
 \frac{|y|\,|R^2y/|y|^2-x|}{R|x-y|}
\)
gives $|\nabla_xG_R(x,y)|\le(\pi|x-y|)^{-1}$, since the reflected
point is at least as far from $x$ as $y$ is. Splitting at
$|y|=|x|/2$ and $|x-y|=|x|/2$ yields
$\int_{\R^2}|x-y|^{-1}(1+|y|)^{-2-\rho}\,dy
\le C_\rho(1+|x|)^{-1}$.
Integration along the segment from $0$ to $x$ then bounds the first
integral by $C_\rho\log(1+|x|)$. These stronger bounds imply both
claims. An $o(1)$ translation changes the weights by uniform factors.
\end{proof}

\begin{lemma}[Signed radial defect]\label{lem:J-negative}
The renormalized integral
\begin{equation}\label{eq:J-def}
 J(\alpha)=\int_{\R^2}
  (U_\alpha+m(\alpha)\log|z|-c_\infty(\alpha))\,dz
 =-\frac14\int_{\R^2}|z|^2F(U_\alpha)\,dz<0
\end{equation}
is finite and continuous on $I$; hence $\sup_IJ<0$.
\end{lemma}
\begin{proof}
For $\mathfrak r_\alpha(r)=U_\alpha(r)+m\log r-c_\infty$,
\begin{equation}\label{eq:R-representation}
 \mathfrak r_\alpha(r)=-\int_r^\infty\frac1s
                  \int_s^\infty tF(U_\alpha(t))\,dt\,ds.
\end{equation}
It is integrable since $\mathfrak r_\alpha=m\log r+O(1)$ near zero and
$\mathfrak r_\alpha=O(r^{2-m})$ at infinity, with $m>4$.
Tonelli's theorem yields
$2\pi\int_0^\infty r\mathfrak r_\alpha(r)\,dr
=-(\pi/2)\int_0^\infty r^3F(U_\alpha(r))\,dr$.
Uniform decay on $I$ and smooth dependence on $\alpha$ give continuity.
\end{proof}

\section{Structure of a non-simple singular cluster}
\label{sec:cluster}

We retain the concentration assumptions of Section~\ref{sec:setting}
and allow both profile regimes.

Fix a non-simple vortex $p_i$, put $n=N_i$, and identify $p_i$ with
the origin in \eqref{eq:isothermal}. Choose $r_{\mathrm{iso}}>0$ so that $B_{4r_{\mathrm{iso}}}$
contains no other vortex or blow-up point. Non-simplicity gives
$x_\eps\in B_{r_{\mathrm{iso}}}$ with
\begin{equation}\label{eq:nonsimple-max}
 \we(x_\eps)+2(n+1)\log|x_\eps|
 =\max_{B_{r_{\mathrm{iso}}}}\{\we(x)+2(n+1)\log|x|\}
 \longrightarrow+\infty.
\end{equation}
The exterior convergence in \eqref{eq:measure-limit} implies
\begin{equation}\label{eq:clustering-scale}
 \de:=|x_\eps|\longrightarrow0.
\end{equation}
Since $u_0(x)=2n\log|x|+O(1)$ and
$\sup_Mu_\eps\le-c_{\mathrm{nt}}<0$,
\eqref{eq:nonsimple-max} also implies
\begin{equation}\label{eq:eta}
 \etae:=\frac{\eps}{\de}\longrightarrow0.
\end{equation}
Set
\begin{equation}\label{eq:first-rescaling-v}
 v_\eps(y):=\we(\de y)-4\pi N_{\tot}\gamma_i(\de y)
 -\Psi_i(p_i+\de y)+2(n+1)\log\de
 +\mathcal K_i^{\comp}(p_i).
\end{equation}
Then
\begin{equation}\label{eq:first-scale-equation-sec3}
 \Delta v_\eps+f_\eps=0,\qquad
 f_\eps:=e^{\varphi_i(\de y)}a_\eps e^{v_\eps}
       (1-\etae^2a_\eps e^{v_\eps}),
\end{equation}
where
\begin{equation}\label{eq:first-scale-coefficient-sec3}
 a_\eps(y):=|y|^{2n}
 \exp\{\mathcal K_i^{\comp}(p_i+\de y)
             -\mathcal K_i^{\comp}(p_i)\}.
\end{equation}
In particular, $\etae^2a_\eps e^{v_\eps}=e^{u_\eps(p_i+\de y)}$.
The density transforms with its area factor:
\begin{equation}\label{eq:first-scale-measure}
 \frac1{2\pi}f_\eps(y)\,dy
 =\frac{F(u_\eps(p_i+\de y))}{2\pi\eps^2}
                     e^{\varphi_i(\de y)}\de^2\,dy.
\end{equation}
Thus this first rescaling preserves normalized local mass.
The smooth background has already been removed in
\eqref{eq:first-rescaling-v}, so \eqref{eq:first-scale-equation-sec3}
has zero right-hand side, including in the surface setting.
We have $f_\eps\ge0$, its integral is uniformly bounded, and
$f_\eps\asymp |y|^{2n}e^{v_\eps}$ on $B_{2r_{\mathrm{iso}}/\de}$. On every fixed
disk disjoint from zero,
\begin{equation}\label{eq:a-uniform}
 C^{-1}\le a_\eps\le C,\qquad
 \|\log a_\eps\|_{C^{4,\sigma}}\le C.
\end{equation}

\begin{lemma}[Regular cores with smooth local coefficients]
\label{lem:local-regular-core}
Let $V_j$ solve, on a fixed disk $B_{4r}(q)$,
\[
 \Delta V_j+\mathcal D_j=0,\qquad
 \mathcal D_j=B_jA_je^{V_j}(1-\tau_j^2A_je^{V_j}),\qquad \tau_j\to0.
\]
Assume $A_j,B_j>0$, $\Delta\log A_j=0$, and, for some $C,c_0>0$,
\[
 C^{-1}\le A_j,B_j\le C,\quad
 \|\log A_j\|_{C^{4,\sigma}}+\|\log B_j\|_{C^{4,\sigma}}\le C,
 \quad \tau_j^2A_je^{V_j}\le1-c_0,\quad
 \int_{B_{4r}(q)}\mathcal D_j\le C.
\]
Suppose $q$ is the only blow-up point in this disk and $V_j\to-\infty$
locally uniformly off $q$. In the Dirichlet Green representation
$V_j=h_j+\int G_{4r}(\cdot,y)\mathcal D_j(y)\,dy$, assume
$\osc_{B_{3r}(q)}h_j\le C$; no bound on its additive constant is required.
Choose $x_j$ as a maximum of $V_j$ on $\overline{B_{3r}(q)}$, and set
\[
 \rho_j=e^{-V_j(x_j)/2},\qquad
 \mathfrak g_j(z)=\rho_j^2\mathcal D_j(x_j+\rho_jz),\qquad
 \widetilde\mu_j=\frac1{2\pi}\int_{B_r(x_j)}\mathcal D_j.
\]
Then $x_j\to q$, $\rho_j\to0$, and, for each $0<\theta<1$,
\begin{equation}\label{eq:regular-core-decay}
 \mathfrak g_j(z)\le C_\theta(1+|z|)^{-4+\theta},
 \qquad |z|\le2r/\rho_j.
\end{equation}
In particular, for $1\le R\le r/\rho_j$,
\begin{equation}\label{eq:regular-core-exhaustion}
 \int_{B_r(x_j)\setminus B_{R\rho_j}(x_j)}\mathcal D_j
 \le C_\theta R^{-2+\theta}.
\end{equation}
Uniformly for $|x-x_j|\le r$,
\begin{equation}\label{eq:local-regular-log}
 V_j(x)=V_j(x_j)-\frac{\widetilde\mu_j}{2}
              \log(1+|x-x_j|^2/\rho_j^2)+O(1).
\end{equation}
After extraction, the whole atom is a single regular profile: its
normalized mass is $4$ in the mean-field case and greater than $4$
in the finite-height case. The constants do not require a positive
lower bound on $\tau_j^2A_j(x_j)e^{V_j(x_j)}$.
\end{lemma}

\begin{proof}
We adapt \cite[Lemmas~2.7, 2.9--2.11]{LinYanAIHP} and
\cite[Theorem~3.1]{Choe}; see also \cite[equation~(3.4)]{CKL}.
Replacing $V_j$ by $V_j+\log A_j$ absorbs the harmonic factor.
On any scale $\ell_j\to0$, $B_j(x_j+\ell_jz)$ converges locally,
after extraction, to a positive constant. The primitive
\[
 \mathcal Q_j(x,V)=B_jA_je^V(1-\tau_j^2A_je^V/2),
 \qquad 0\le\mathcal Q_j\le C\mathcal D_j
\]
has explicit spatial derivative
$\mathcal D_j\nabla\log A_j+\mathcal Q_j\nabla\log B_j$.
The rescaled translation identities therefore have coefficient error
at most $C\ell_j\int\mathcal D_j$; the bounded harmonic oscillation
gives a further $O(\ell_j)$ gradient error. Sources outside a rescaled
disk of large radius $T$ contribute at most
$CT^{-1}\int\mathcal D_j$ on fixed compact sets.
Letting $j\to\infty$ and then $T\to\infty$ gives the constant-coefficient
profile equations and interaction identities, also for smooth metric
factors $B_j$, without using core exhaustion.

At the maximum scale $\rho_j$, local compactness gives an entire
regular profile. After constant normalizations, its type is determined
by the limit of $\tau_j^2A_j(x_j)e^{V_j(x_j)}$: zero gives a Liouville
profile of normalized mass $4$, and a positive limit gives a
Chern--Simons profile of mass greater than $4$;
see \cite[pp.~840--841]{CKL} and \cite[Lemmas~2.5--2.6]{LinYanAIHP}.
Thus a sufficiently large fixed rescaled disk has mass at least
$4-\theta$ for large $j$.

There can be no positive mass in a disk centered a distance
$\ell_j\gg\rho_j$ from $x_j$ and of radius $\ell_j/2$.
Otherwise, after extraction, $\ell_j\to0$, since disks remaining
a fixed distance from $q$ have vanishing mass. Set
\(W_j(z)=V_j(x_j+\ell_jz)+2\log\ell_j.\)
Since $\tau_j=O(\rho_j)$, we have $\lambda_j:=\tau_j/\ell_j\to0$.
Writing $A_j^\ell(z)=A_j(x_j+\ell_jz)$ and similarly for $B_j^\ell$, we have
\[
 \Delta W_j+C_je^{W_j}=0,\qquad
 C_j=B_j^\ell A_j^\ell(1-\lambda_j^2A_j^\ell e^{W_j}),
 \qquad 0<c\le C_j\le C.
\]
Thus $\int e^{W_j}\le C$ on each fixed rescaled disk, while
$W_j(0)=2\log(\ell_j/\rho_j)\to+\infty$.
The concentration alternative of \cite[Theorem~3]{BM}, applied
diagonally on expanding fixed rescaled disks, gives atomic convergence
and decay away from a finite set: the atom threshold and total mass
bound ensure finiteness. The original core supplies an atom at zero,
and the positive mass at unit distance forces another, without using
core exhaustion. Write $\xi_a$ and $\mu_a>0$ for the atom locations
and normalized masses. Apply the translation identity on a small
circle enclosing only $\xi_a$, first letting $j\to\infty$ and then
shrinking the circle. Its boundary terms use $W_j\to-\infty$ and
the exterior Green-gradient limit; no decomposition or exhaustion
of the enclosed atom is needed. The identity gives
\[
 \sum_{b\ne a}\mu_b
        \frac{\xi_a-\xi_b}{|\xi_a-\xi_b|^2}=0.
\]
The coefficient, harmonic, and remote-source errors vanish in the
order of limits specified above.
Taking scalar products with $\mu_a\xi_a$ and summing gives
$\sum_{a<b}\mu_a\mu_b=0$, a contradiction. Disks remaining a fixed
distance from $q$ already have vanishing mass. This argument yields
the following \emph{moving-disk} estimate, rather than a conclusion
about the mass of the entire neck:
\begin{equation}\label{eq:regular-moving-disk}
 \lim_{R\to\infty}\limsup_{j\to\infty}
 \sup_{R\rho_j\le |x-x_j|\le2r}
 \int_{B_{|x-x_j|/2}(x)}\mathcal D_j=0.
\end{equation}
The disks lie in $B_{4r}(q)$ for large $j$. This is the local form of
\cite[Lemma~2.7]{LinYanAIHP}, justified by the coefficient checks above.

We now obtain a summable annular bound from this local estimate.
Put $\widehat V_j(z)=V_j(x_j+\rho_jz)-V_j(x_j)$ and
$\Omega_j=(B_{4r}(q)-x_j)/\rho_j$. Subtracting the Dirichlet Green
formula at $x_j$ gives, uniformly for $|z|\le2r/\rho_j$,
\begin{equation}\label{eq:regular-green-difference}
 \widehat V_j(z)=\frac1{2\pi}\int_{\Omega_j}
      \log\frac{|y|}{|z-y|}\,\mathfrak g_j(y)\,dy+O(1).
\end{equation}
The harmonic oscillation, regular Green part, and total mass control
the $O(1)$ term uniformly.
The maximum normalization implies $\mathfrak g_j\le C$ on every
rescaled subset corresponding to $B_{3r}(q)$.

Fix $0<\theta<1$. Choose a fixed $R_\theta$ so that, for large $j$,
\(
 \frac1{2\pi}\int_{B_{R_\theta}}\mathfrak g_j
       \ge4-\theta/4.
\)
This follows from local convergence to the first entire profile,
without any assertion about its neck. For $s=|z|\ge2R_\theta$,
decompose the Green integral into the three disjoint regions
\[
 \begin{split}
 \mathcal E_{1,j}&=B_{R_\theta},\quad
 \mathcal E_{2,j}=(\Omega_j\setminus B_{R_\theta})\cap B_{s/2}(z),\quad
 \mathcal E_{3,j}=\Omega_j\setminus
                         (\mathcal E_{1,j}\cup\mathcal E_{2,j}).
 \end{split}
\]
The first region contributes at most
$-(4-\theta/4)\log s+C_\theta$. The fixed loss $\theta/4$ is retained
before multiplication by the unbounded factor $\log s$.
On $\mathcal E_{3,j}$ the kernel is bounded above by $\log3$,
since $|y|\le|y-z|+s\le3|y-z|$. The total mass controls this term.
The remaining region contains the moving logarithmic singularity;
we use \eqref{eq:regular-moving-disk} there.

For the remaining disk choose $0<a<1$. The local bound on
$\mathfrak g_j$ gives
\[
 \frac1{2\pi}\int_{B_a(z)}
       \log\frac{|y|}{|z-y|}\,\mathfrak g_j(y)\,dy
       \le Ca^2\log(2+s)+C_a.
\]
On $B_{s/2}(z)\setminus B_a(z)$, the same integral is at most
\(
 \frac1{2\pi}\left(\int_{B_{s/2}(z)}\mathfrak g_j\right)
                   \log\frac{Cs}{a}.
\)
First choose $a$ so that $Ca^2<\theta/4$, and then use
\eqref{eq:regular-moving-disk} to make the disk mass divided by
$2\pi$ less than $\theta/4$, uniformly for sufficiently large $s$
and $j$. Bounded $s$ are covered by the maximum normalization.
Consequently
\[
 \widehat V_j(z)\le-(4-\theta)\log(1+|z|)+C_\theta,
 \qquad |z|\le2r/\rho_j.
\]
The coefficient bounds imply $\mathfrak g_j\le C e^{\widehat V_j}$,
which proves \eqref{eq:regular-core-decay}.

For clarity, the passage to the whole neck is a geometric sum.
For $R\ge1$ put
\[
 A_{j,k}=\{2^kR\rho_j\le |x-x_j|<2^{k+1}R\rho_j\}
                 \cap B_{2r}(x_j),\qquad k\ge0.
\]
The decay just proved gives
\begin{equation}\label{eq:regular-dyadic-neck}
 \int_{A_{j,k}}\mathcal D_j
       \le C_\theta(2^kR)^{-2+\theta},\qquad
 \int_{B_{2r}(x_j)\setminus B_{R\rho_j}(x_j)}\mathcal D_j
       \le \frac{C_\theta R^{-2+\theta}}{1-2^{-2+\theta}}.
\end{equation}
The bound is independent of the number of nonempty annuli, so it
excludes diffuse neck mass and proves \eqref{eq:regular-core-exhaustion}
in both regimes. Exhaustion outside the whole cluster is proved
separately in Proposition~\ref{prop:polygon}.

Finally, \eqref{eq:regular-core-decay} gives the uniform logarithmic
moment bound
\[
 \int_{|y|\le2r/\rho_j}\log(2+|y|)\,\mathfrak g_j(y)\,dy
 \le C_\theta\int_0^\infty t(1+t)^{-4+\theta}\log(2+t)\,dt<\infty.
\]
Together with local integrability of the logarithmic kernel, this
allows the Green-kernel decomposition in
\cite[proof of Theorem~3.1]{Choe} to bound both logarithmic remainders.
The leading coefficient is the actual finite-disk mass,
as in \cite[equation~(32)]{Choe}. Replacing the mass
of $B_{2r}(x_j)$ by that of $B_r(x_j)$ costs
$O(\rho_j^{2-\theta}|\log\rho_j|)=o(1)$.
This gives \eqref{eq:local-regular-log}; the bounded difference
between $\log(1+t)$ and $\tfrac12\log(1+t^2)$ is immaterial.
\end{proof}

\begin{lemma}[Compact-scale comparison and balance]
\label{lem:compact-cluster-balance}
After passage to a subsequence there are finitely many distinct
nonzero centers $q_t$, $t=0,\ldots,S-1$, with $|q_0|=1$, and
points $q_{t,\eps}\to q_t$ chosen as maxima on fixed disjoint small
disks about $q_t$. Write $b_t>0$ for their limiting masses
and $c$ for a possible mass at zero. Then $c=0$, $S\ge2$, and
\begin{equation}\label{eq:weighted-location-system}
 2n\frac{q_t}{|q_t|^2}
 -\sum_{s\ne t}b_s\frac{q_t-q_s}{|q_t-q_s|^2}=0.
\end{equation}
Moreover, $v_\eps\to-\infty$ uniformly on a small disk about zero.
\end{lemma}

\begin{proof}
Apply the concentration alternative of \cite{ChoeKim} and
\cite[Section~2]{LinYanAIHP} on connected annuli containing the
unit-circle blow-up point in \eqref{eq:nonsimple-max}, and exhaust
$\R^2\setminus\{0\}$ diagonally. Here $f_\eps\asymp e^{v_\eps}$
with bounded source mass. On each isolated atom disk, the localized
Green formula has bounded harmonic oscillation: omitted sources have
bounded mass and stay a fixed distance away, while rescaled smooth
Green terms have bounded derivatives. Its additive constant is irrelevant.
Lemma~\ref{lem:local-regular-core} applies with
\(A_j=a_\eps,\quad B_j=e^{\varphi_i(\de\,\cdot)},\quad\tau_j=\etae\),
since $\log a_\eps$ is harmonic away from zero,
\eqref{eq:a-uniform} bounds the coefficients, and
$1-\etae^2a_\eps e^{v_\eps}\ge1-e^{-c_{\mathrm{nt}}}>0$.
Thus every regular atom is simple with mass $b_t\ge4$, and the
total mass bound makes their number finite.

The Green representation gives, on compact subsets of
$\R^2\setminus\{0,q_0,\ldots,q_{S-1}\}$,
\begin{equation}\label{eq:compact-green-gradient}
 \nabla v_\eps(y)\longrightarrow
 -c\frac{y}{|y|^2}
 -\sum_t b_t\frac{y-q_t}{|y-q_t|^2}.
\end{equation}
Indeed, for a fixed large $R$, convergence of the local source
measures and core centers, together with the locally uniform decay
of $f_\eps$ away from the atoms, gives an $o_\eps(1)$ error on such
compact sets. The rescaled smooth Green terms have gradient $O(\de)$,
whereas the sources outside $B_R$ contribute at most $C/R$ times
the total mass. Thus the remainder is
$o_\eps(1)+O(\de)+O(R^{-1})$; no rate is asserted for the local
measure convergence. Letting first $\eps\to0$ and then $R\to\infty$
proves \eqref{eq:compact-green-gradient}; subtraction at two points
also gives bounded exterior oscillation.

To pass to the radial Pohozaev identity, first fix a small radius
whose boundary contains no regular atom, let $\eps\to0$, and then
let the radius tend to zero. The nonnegative densities
$e^{\varphi_i(\de y)}\etae^2a_\eps^2e^{2v_\eps}$ are bounded by
$C f_\eps$, so their measures admit a weakly convergent subsequence.
Using continuity radii for the limiting measure, the identity becomes
\[
 \pi c^2=4\pi(n+1)c+
 \lim_{r\downarrow0}\lim_{\eps\to0}
 \int_{B_r}e^{\varphi_i(\de y)}
                  \etae^2a_\eps^2e^{2v_\eps}\,dy.
\]
It follows that $c=0$ or $c\ge4(n+1)$. To evaluate the translational
identity, use the primitive
\(\mathcal P_\eps(y,v)=e^{\varphi_i(\de y)}a_\eps e^v (1-\etae^2a_\eps e^v/2).\)
Its explicit spatial derivative is
$f_\eps\nabla\log a_\eps+\mathcal P_\eps\nabla\varphi_i(\de y)$.
The second term tends to zero on compact sets, since
$|\mathcal P_\eps|\le C f_\eps$ and $\nabla_y\varphi_i(\de y)\to0$.
Together with \eqref{eq:compact-green-gradient}, the identity on
each core disk therefore gives
\(
 (2n-c)\frac{q_t}{|q_t|^2}
 -\sum_{s\ne t}b_s\frac{q_t-q_s}{|q_t-q_s|^2}=0.
\)
Taking scalar products with $b_tq_t$ and summing yields
\begin{equation}\label{eq:weighted-cluster-virial}
 (2n-c)\sum_t b_t=\sum_{s<t}b_sb_t.
\end{equation}
Its right-hand side is nonnegative and $\sum_t b_t>0$, so
$c\le2n$. This contradicts $c\ge4(n+1)$ whenever $c>0$.
Hence $c=0$, without any equality assumption on the regular masses.
Since $n>0$, \eqref{eq:weighted-cluster-virial} also excludes $S=1$.

Finally, fix a disk $B_\varrho(0)$ containing no regular atom.
Its boundary values tend uniformly to $-\infty$ and, since $c=0$,
$\int_{B_\varrho}f_\eps\to0$. Write
$v_\eps=\mathfrak h_\eps+\mathfrak w_\eps$, where
$\mathfrak h_\eps$ is the harmonic extension of the boundary data
and $\mathfrak w_\eps$ is the zero-boundary Green potential of
$f_\eps$. The maximum principle gives
$\sup\mathfrak h_\eps\to-\infty$ and $\mathfrak w_\eps\ge0$.
For any fixed $p>1$, the small-mass exponential estimate of
\cite{BM} gives
\[
 \int_{B_\varrho}e^{p\mathfrak w_\eps}\le C_p,\qquad
 \|f_\eps\|_{L^p(B_\varrho)}
   \le C e^{\sup\mathfrak h_\eps}
                      \|e^{\mathfrak w_\eps}\|_{L^p(B_\varrho)}
   \longrightarrow0.
\]
Here $f_\eps\le C e^{v_\eps}$ because $n>0$ and the remaining
coefficient is bounded on this fixed disk. The Dirichlet
$W^{2,p}$ estimate and its embedding into $C^0$ give
$\|\mathfrak w_\eps\|_\infty\to0$. Therefore
$v_\eps\to-\infty$ uniformly on a smaller disk about zero.
\end{proof}

The lemma also applies after replacing $\de$ by any scale
$s_\eps\to0$ with $\eps/s_\eps\to0$ on which a nonzero regular
concentration point occurs. Its proof allows arbitrary positive regular masses.

Choose $r>0$ so that the disks $B_{2r}(q_{t,\eps})$ are disjoint
and avoid zero, and define
\begin{equation}\label{eq:first-local-mass}
 \widetilde m_{t,\eps}:=\frac1{2\pi}
       \int_{B_r(q_{t,\eps})}f_\eps,\qquad
 \vartheta_{t,\eps}:=v_\eps(q_{t,\eps})+2\log\etae.
\end{equation}

The core width in the first-scale coordinates and its counterpart
in the original coordinates are, respectively,
\[
 \rho_{t,\eps}:=e^{-v_\eps(q_{t,\eps})/2}
       =\etae e^{-\vartheta_{t,\eps}/2},
 \qquad
 \de\rho_{t,\eps}=\eps e^{-\vartheta_{t,\eps}/2}.
\]
These widths are distinct from the cluster scale $\de$.
The normalized physical height is
$\vartheta_{t,\eps}+\log a_\eps(q_{t,\eps})
=u_\eps(p_i+\de q_{t,\eps})$.

\begin{lemma}[Uniform logarithmic comparison]\label{lem:isolated-core-log}
For each selected core, uniformly on $B_r(q_{t,\eps})$,
\begin{equation}\label{eq:isolated-core-log}
 v_\eps(y)=v_\eps(q_{t,\eps})
 -\frac{\widetilde m_{t,\eps}}2
 \log\bigl(1+e^{v_\eps(q_{t,\eps})}|y-q_{t,\eps}|^2\bigr)+O(1).
\end{equation}
The constant is uniform in $t,\eps$ and in both profile regimes.
\end{lemma}
\begin{proof}
Apply \eqref{eq:local-regular-log} with
$\rho_j=\exp(-v_\eps(q_{t,\eps})/2)$ and mass
\eqref{eq:first-local-mass}. Lemma~\ref{lem:compact-cluster-balance}
verifies the hypotheses; choose the maxima on slightly larger
isolated-core disks.
\end{proof}

\begin{lemma}[Equal masses and a common profile type]
\label{lem:equal-core-masses}
There is $m\ge4$ such that
\begin{equation}\label{eq:common-local-mass}
 \widetilde m_{t,\eps}\longrightarrow m
 \quad(t=0,\ldots,S-1).
\end{equation}
Either all cores are of mean-field type and $m=4$, or all are of
Chern--Simons type and $m>4$.
\end{lemma}

\begin{proof}
As in \cite[Lemma~3.3]{CKL} and
\cite[proof of Theorem~1.1]{LinYanAIHP}, compare
\eqref{eq:isolated-core-log} at fixed exterior distances.
The bounded exterior oscillation and cancellation of the common
term $-2\log\etae$ give
\begin{equation}\label{eq:qualitative-mass-comparison}
 \widetilde m_{t,\eps}\log\etae
 -\frac{\widetilde m_{t,\eps}-2}{2}\vartheta_{t,\eps}
 =\widetilde m_{s,\eps}\log\etae
 -\frac{\widetilde m_{s,\eps}-2}{2}\vartheta_{s,\eps}+O(1).
\end{equation}
After extraction, each $\vartheta_{t,\eps}$ is bounded or tends to
$-\infty$. Regular profile classification gives mass $>4$ in the
first case and mass $4$ in the second. If the $t$-th height is
bounded and the $s$-th tends to $-\infty$, the last identity reads
\[
 (\widetilde m_{t,\eps}-\widetilde m_{s,\eps})\log\etae
 =-\frac{\widetilde m_{s,\eps}-2}{2}\vartheta_{s,\eps}+O(1),
\]
whose two sides tend to opposite infinities. Hence the types cannot
coexist. For bounded heights, division by $\log\etae$ proves
equality of the limiting masses; in the other case every mass is
$4$. In the Chern--Simons case, strict monotonicity of the radial
mass map identifies a common limiting normalized physical height
$\vartheta_{t,\eps}+\log a_\eps(q_{t,\eps})$.
\end{proof}

\begin{proposition}[Mass and polygon identities]\label{prop:polygon}
After rotation and relabeling,
\begin{equation}\label{eq:regular-polygon}
 q_t=e_t:=e^{2\pi\mathrm i t/S},\qquad t=0,\ldots,S-1.
\end{equation}
The common core mass and the full mass at the vortex satisfy
\begin{align}
 m(S-1)&=4N_i,\label{eq:mass-location}\\
 M_{p_i}&=Sm=\frac{4N_iS}{S-1}.\label{eq:local-total-mass}
\end{align}
In particular,
\begin{equation}\label{eq:S-bound}
 1\le S-1\le N_i,
\end{equation}
with equality on the right exactly in the mean-field case.
\end{proposition}

\begin{proof}
The balance equation becomes
\begin{equation}\label{eq:location-system}
 2n\frac{q_t}{|q_t|^2}
 -m\sum_{s\ne t}\frac{q_t-q_s}{|q_t-q_s|^2}=0.
\end{equation}
Taking scalar products with $q_t$ and pairing the interactions
gives $2nS=mS(S-1)/2$, proving \eqref{eq:mass-location}.
Identify the plane with $\mathbb C$ and rotate so that $q_0=1$.
Complex conjugation of \eqref{eq:location-system} gives
\begin{equation}\label{eq:complex-location}
 \sum_{s\ne t}\frac{q_t}{q_t-q_s}=\frac{S-1}{2}.
\end{equation}
For $\mathfrak p(z)=\prod_t(z-q_t)$, the relation
$\mathfrak p''(q_t)/(2\mathfrak p'(q_t))=\sum_{s\ne t}(q_t-q_s)^{-1}$ implies that
$z\mathfrak p''-(S-1)\mathfrak p'$ vanishes at all $S$ distinct roots. Its degree is
less than $S$, so it vanishes identically. Thus $\mathfrak p(z)=z^S-1$,
which proves \eqref{eq:regular-polygon}.

It remains to identify the mass at the original vortex, including
possible larger scales. Fix $R_0>2$ containing the finite rescaled
centers. We first claim
\begin{equation}\label{eq:exterior-spherical-bound}
 v_\eps(y)+2(n+1)\log|y|\le C,
 \qquad R_0\le|y|\le2r_{\mathrm{iso}}/\de.
\end{equation}
Otherwise there are $y_\eps$, $T_\eps=|y_\eps|\to\infty$, with
$\de T_\eps\to0$ and
$v_\eps(y_\eps)+2(n+1)\log T_\eps\to\infty$; the two scale
limits follow from compact-scale and original-scale exterior
convergence. The further rescaling
\(g_\eps(z)=v_\eps(T_\eps z)+2(n+1)\log T_\eps\)
is taken at the physical scale $\ell_\eps=\de T_\eps\to0$.
To keep all coefficient factors explicit, put
\[
 \widehat\eta_\eps=\frac{\etae}{T_\eps}=\frac{\eps}{\ell_\eps},
 \qquad
 \widehat a_\eps(z)=T_\eps^{-2n}a_\eps(T_\eps z)
 =|z|^{2n}\exp\{\mathcal K_i^{\comp}(p_i+\ell_\eps z)
                         -\mathcal K_i^{\comp}(p_i)\}.
\]
The exact rescaled equation and density are
\begin{equation}\label{eq:exterior-rescaled-equation}
 \begin{split}
 \Delta g_\eps+\widehat f_\eps&=0,\\
 \widehat f_\eps(z):=T_\eps^2f_\eps(T_\eps z)
 &=e^{\varphi_i(\ell_\eps z)}\widehat a_\eps(z)e^{g_\eps(z)}
           (1-\widehat\eta_\eps^2\widehat a_\eps(z)e^{g_\eps(z)}).
 \end{split}
\end{equation}
In particular $\widehat\eta_\eps\to0$, the negative upper bound
on the physical solution is preserved, and the total density
mass is unchanged by this rescaling. The assumed failure yields
a regular concentration point on the unit circle. The old
cluster collapses to zero and supplies a normalized central mass
\[
 \widehat c:=\lim_{\varrho\downarrow0}\lim_{\eps\to0}
       \frac1{2\pi}\int_{B_\varrho}\widehat f_\eps
       \ge\sum_t b_t=Sm>0.
\]
Indeed, for fixed $\varrho>0$ the corresponding disk
$B_{\varrho T_\eps}$ eventually contains every fixed disk
surrounding the old cluster. The local coefficient hypotheses of
Lemma~\ref{lem:compact-cluster-balance} hold with $\de$ replaced
by $\ell_\eps$, as noted after its proof. That lemma forces
$\widehat c=0$, a contradiction. This uses its weighted
central-mass exclusion, not an energy identity for the old
cluster at all larger scales. It proves
\eqref{eq:exterior-spherical-bound}.

We now exclude diffuse neck mass as well. The global Green formula,
localized on $B_{2r_{\mathrm{iso}}}$ and evaluated on $B_{r_{\mathrm{iso}}}$, gives
\begin{equation}\label{eq:exterior-green-difference}
 v_\eps(y)-v_\eps(0)
 =\frac1{2\pi}\int_{B_{2r_{\mathrm{iso}}/\de}}
       \log\frac{|z|}{|y-z|}f_\eps(z)\,dz+O(1),
 \qquad |y|\le r_{\mathrm{iso}}/\de.
\end{equation}
The remainder is uniform: the regular Green part and the fixed
regularizations are smooth, and the omitted sources are a fixed
positive distance from the evaluation disk. For $|y|\ge2R_0$,
split the integral into $|z|\le R_0$, the set
$|z|>R_0$, $|z-y|<|y|/2$, and its complement. The first term is
\[
 -\widetilde M_\eps(R_0)\log|y|+O(1),\qquad
 \widetilde M_\eps(R_0):=\frac1{2\pi}\int_{B_{R_0}}f_\eps
                          \longrightarrow Sm.
\]
Its logarithmic moment at zero is bounded because
$v_\eps\to-\infty$ uniformly near zero. On the second set,
\eqref{eq:exterior-spherical-bound} gives $f_\eps(z)\le C|z|^{-2}$;
rescaling by $|y|$ bounds the integrable logarithmic singularity.
On the last set, $|z|\le3|z-y|$, and the upper bound follows from
the total mass. Consequently, for each fixed $\kappa>0$ and all
sufficiently small $\eps$,
\begin{equation}\label{eq:exterior-fixed-loss}
 v_\eps(y)\le v_\eps(0)-(Sm-\kappa)\log|y|+C_\kappa,
 \qquad 2R_0\le|y|\le r_{\mathrm{iso}}/\de.
\end{equation}
Here the fixed loss is retained before multiplying by $\log|y|$;
no rate for $\widetilde M_\eps(R_0)\to Sm$ is needed.
Since $Sm=4n+m$ and $m\ge4$, choose
$0<\kappa<Sm-2n-2$. Then
\begin{equation}\label{eq:complete-cluster-mass}
 \int_{B_{r_{\mathrm{iso}}/\de}\setminus B_{2R_0}}f_\eps
 \le C_\kappa e^{v_\eps(0)}
       \int_{2R_0}^\infty t^{2n-Sm+\kappa+1}\,dt
 \longrightarrow0.
\end{equation}
On the fixed inner disk, the measure convergence gives
\(\frac1{2\pi}\int_{B_{2R_0}}f_\eps=Sm+o_\eps(1).\)
Combining this with \eqref{eq:complete-cluster-mass} and then
returning to the original coordinates proves
\eqref{eq:local-total-mass}. Finally, \eqref{eq:S-bound} follows
from $m\ge4$, with equality precisely when $m=4$.
\end{proof}

\begin{remark}
The argument applies to every real strength $n>0$ and to the
surface setting above. Away from zero the weight is smooth; at
zero the radial computation uses
$y\cdot\nabla\log|y|^{2n}=2n$ and follows by integration on
punctured disks. Smooth metric terms disappear under rescaling.
In particular, mean-field clusters require $S=n+1$, hence an
integer strength $n$, whereas finite-height clusters have $S<n+1$.
\end{remark}

\begin{proof}[Proof of Theorem~\ref{thm:ratio-free-four-unit-gap}]
For a non-simple cluster at $p_i$, Proposition~\ref{prop:polygon} gives
$k=S-1\in\mathscr A_i$ and $M_{p_i}=M_{i,k}$.
Every additional blow-up point has mass at least $4$ by
\eqref{eq:measure-limit}, while the total mass is $2N_{\tot}$ by
\eqref{eq:global-mass}. Hence
\[
 2N_{\tot}-M_{i,k}=0\quad\hbox{or}\quad
 2N_{\tot}-M_{i,k}\ge4.
\]
This is \eqref{eq:necessary-alternative} and contradicts
\eqref{eq:ratio-free-four-unit-gap}. Hence all singular blow-up
points are simple. Near a regular blow-up point, choose a vortex-free
conformal chart and a smooth $\Gamma$ with $\Delta\Gamma=e^\varphi/|M|_g$.
Set
\[
 V_\eps=w_\eps-4\pi N_{\tot}\Gamma,\qquad
 A=e^{u_0+4\pi N_{\tot}\Gamma},\qquad B=e^\varphi.
\]
Then $\Delta\log A=0$, and $V_\eps$ satisfies the equation of
Lemma~\ref{lem:local-regular-core} with $\tau=\eps$.
The global Green formula bounds the harmonic oscillation.
Adding back the fixed smooth term in its logarithmic estimate
gives regular simplicity at a nearby local maximum of $w_\eps$.
\end{proof}

\section{Quantitative estimates for finite-height cores}
\label{sec:pointwise}

We quantify finite-height cores in a collapsing non-simple cluster.
The proof of Theorem~\ref{prop:refined-profile} and its moment-sampling
consequences are in Appendix~\ref{app:pointwise-proofs}.

Under Assumption~\ref{ass:profile}, use the first-scale equation,
centers $q_{t,\eps}$, and actual masses $\widetilde m_{t,\eps}$ from
Section~\ref{sec:cluster}. Choose a fixed radius $r>0$ such that
the isolated-core disks $B_r(q_{t,\eps})$ are pairwise disjoint. Let $\widetilde{\mathfrak h}_{t,\eps}$ be the harmonic
extension of $v_\eps-\fint_{\partial B_r(q_{t,\eps})}v_\eps\,dS$.
Let $q_{t,\eps}^*$ maximize
$v_\eps-\widetilde{\mathfrak h}_{t,\eps}$ on the disk, and put
\[
 \mathfrak h_{t,\eps}
 =\widetilde{\mathfrak h}_{t,\eps}
  -\widetilde{\mathfrak h}_{t,\eps}(q_{t,\eps}^*),
 \qquad
 \mathfrak u_{t,\eps}=v_\eps-\mathfrak h_{t,\eps}.
\]
Thus $\mathfrak u_{t,\eps}$ has constant boundary trace and
$\mathfrak h_{t,\eps}(q_{t,\eps}^*)=0$. Define
\[
 \begin{aligned}
 H_{t,\eps}=a_\eps e^{\mathfrak h_{t,\eps}},
 \ \ s_{t,\eps}=\etae e^{-\varphi_i(\de q_{t,\eps}^*)/2},\ \  
 \alpha_{t,\eps}
 =v_\eps(q_{t,\eps}^*)+2\log\etae
                         +\log a_\eps(q_{t,\eps}^*),
 \ \ m_{t,\eps}=m(\alpha_{t,\eps}),
 \end{aligned}
\]
and
\[
 \begin{aligned}
 \widehat v_{t,\eps}(z)
  =\mathfrak u_{t,\eps}(q_{t,\eps}^*+s_{t,\eps}z)
             +2\log\etae+\log H_{t,\eps}(q_{t,\eps}^*),\ \
 U_{t,\eps} =U_{\alpha_{t,\eps}},\ \
 W_{t,\eps}=\widehat v_{t,\eps}-U_{t,\eps}.
 \end{aligned}
\]
Here $m_*>4$ is the uniform lower bound chosen in
\eqref{eq:profile-hyp}. Appendix~\ref{app:core-normalization} verifies
these normalizations and the preliminary estimates.

\begin{theorem}[Pointwise approximation at each Chern--Simons bubble]
\label{prop:refined-profile}
Assume the finite-height cluster setting and the local normalizations above.
Fix
\begin{equation}
  0<\sigma<\min\{1,m_*-4\}.
  \label{eq:sigma-choice}
\end{equation}
Uniformly for $t=0,\dots,S-1$, the following assertions hold.

\smallskip
\noindent\textup{(i)}   
For $|z|\le r/(2s_{t,\eps})$, we have
\begin{equation}\label{eq:global-pointwise-O2}
 |W_{t,\eps}(z)|+|\nabla W_{t,\eps}(z)|
 \le C\bigl(\etae^2+\etae\de^2\bigr).
\end{equation}
If $\varphi_i$ is constant on the local chart, the term
$\etae\de^2$ in \eqref{eq:global-pointwise-O2} is absent.

\smallskip
\noindent\textup{(ii)} \(
  \left|
  \nabla\log H_{t,\eps}(q_{t,\eps}^*)
  +\frac{m_{t,\eps}}4
   \de\nabla_y\varphi_i(y)\Big|_{y=\de q_{t,\eps}^*}
  \right|
  \le C\etae^2.
\)

\smallskip
\noindent\textup{(iii)}
Writing
$\mathbf{d}_{t,\eps}=q_{t,\eps}^*-q_{t,\eps}$ and
$\kappa_{t,\eps}=F(\alpha_{t,\eps})$, one has
\(
  \mathbf{d}_{t,\eps}
  =-\frac{2s_{t,\eps}^2}{\kappa_{t,\eps}}
  \nabla\mathfrak h_{t,\eps}(q_{t,\eps}^*)
  +O(\etae^3+\etae^4\de^2).
  \)
In particular, $|q_{t,\eps}^*-q_{t,\eps}|\le C\etae^2$.

\smallskip
\noindent\textup{(iv)}
Let
\begin{equation}\label{eq:B-tracefree}
\begin{aligned}
 \mathsf B_{t,\eps}
   &:=D^2\log H_{t,\eps}(q_{t,\eps}^*),
 &\mathsf B_{t,\eps}^{\circ}
   &:=\mathsf B_{t,\eps}-\tfrac12(\tr\mathsf B_{t,\eps})I_2,\\
 \mathsf C_{t,\eps}
   &:=\de^2D^2_y\varphi_i(y)\Big|_{y=\de q^*_{t,\eps}},
 &\mathsf C_{t,\eps}^{\circ}
   &:=\mathsf C_{t,\eps}-\tfrac12(\tr\mathsf C_{t,\eps})I_2.
\end{aligned}
\end{equation}
There is a unique bounded pure second Fourier mode $\Psi_{2,t,\eps}$
satisfying
\begin{equation}\label{eq:Psi2-equation}
\left\{
\begin{aligned}
 L_{\alpha_{t,\eps}}\Psi_{2,t,\eps}
   &=-\tfrac12(z^T\mathsf B_{t,\eps}^{\circ}z)F'(U_{t,\eps}(z))
     -\tfrac12(z^T\mathsf C_{t,\eps}^{\circ}z)F(U_{t,\eps}(z)),\\
 \Psi_{2,t,\eps}(0)&=0,\qquad \nabla\Psi_{2,t,\eps}(0)=0.
\end{aligned}
\right.
\end{equation}
With the conformal first Fourier correction $\Xi_{1,t,\eps}$
defined in \eqref{eq:Xi1-definition} of
Appendix~\ref{app:pointwise-proof},
for $|z|\le 3r/(4s_{t,\eps})$ we have
\begin{equation}\label{eq:second-mode-refinement}
 \left|\widehat v_{t,\eps}(z)-U_{t,\eps}(z)
       -\Xi_{1,t,\eps}(z)-s_{t,\eps}^2\Psi_{2,t,\eps}(z)\right|
 \le C\bigl(\etae^{2+\sigma}+\etae^2\de^2\bigr)(1+|z|)^\sigma.
\end{equation}

\smallskip
\noindent\textup{(v)} We have
\begin{equation}
  \widetilde m_{t,\eps}=m_{t,\eps}
  +O(\etae^{2+\sigma}+\etae^2\de^2).
  \label{eq:local-mass-refined}
\end{equation}
\end{theorem}

The effective gradient and $\Xi_{1,t,\eps}$ account for the surface
geometry, with $\|\Xi_{1,t,\eps}\|_{C^1}\le C\etae\de^2$ on the
rescaled core disk. The periodic construction in
Section~\ref{sec:cs-construction} derives its moments directly from
the projected solution.

\section{Construction of a finite-height Chern--Simons cluster}
\label{sec:cs-construction}

We prove Theorem~\ref{cs:thm:existence} through a projected periodic
problem, independently of Theorem~\ref{prop:refined-profile}.
Lemma~\ref{cs:lem:moments} derives the required signed moments directly
from the projected solution, including the decisive second angular moment.
Isotropic moments cancel against harmonic fields as in
Proposition~\ref{prop:moment-sampling}.

\subsection{Normalization for the periodic construction}

Use the torus, points $a,e$, Green function $G$, and background
$u_0$ of Theorem~\ref{cs:thm:existence}; see
\eqref{cs:eq:green}--\eqref{cs:eq:u0}. Locally,
\[
 G(x,0)=-\frac1{2\pi}\log|x|+R(x),\qquad \Delta R=1,
 \qquad \Delta u_0=8\pi\delta_0+24\pi\delta_a-32\pi.
\]
The profile in the theorem is $U=U_{\alpha_8}$, where
$m(\alpha_8)=8$ in the notation of Section~\ref{sec:setting}.
Define $\Phi$ and $\mathfrak b>0$ by Lemma~\ref{cs:lem:profile} below, and put
\begin{equation}\label{cs:eq:lambda}
 \lambda=24\pi\bigl(G_{xx}(a,0)-\tfrac12\bigr),\qquad
 d_*=\left(\frac{\mathfrak b}{2\lambda}\right)^{1/6}.
\end{equation}

We use the periodic matching and projected fixed-point scheme of
\cite[Section~3, especially Proposition~3.3]{LinYanAIHP},
establishing below the inverse and interaction estimates uniformly
as the centers approach one another.

\subsection{Two properties of the entire profile}

The radial family and $m'(\alpha_8)>0$ are recalled in
Section~\ref{sec:setting}; see \cite[pp.~840--841 and Lemma~2.4]{CKL}.
With $c_8=c_\infty(\alpha_8)$, the tail estimates of
Section~\ref{sec:linearized} specialize to
\begin{equation}\label{cs:eq:tails}
 U(r)=-8\log r+c_8+O(r^{-6}),\quad
 U'(r)=-8r^{-1}+O(r^{-7}),\quad
 |F^{(j)}(U(r))|\le C_j(1+r)^{-8}.
\end{equation}
Here the last inequality is used for $j\ge0$ in a fixed finite range.

\begin{lemma}\label{cs:lem:profile}
The bounded kernel of $\Delta+F'(U)$ on $\R^2$ is
$\operatorname{span}\{\partial_1U,\partial_2U\}$. The radial function
$Z_0=\left.\partial_\alpha U_\alpha\right|_{\alpha=\alpha_8}/m'(\alpha_8)$ satisfies
\begin{equation}\label{cs:eq:radial-kernel}
 Z_0(r)=-\log r+O(1),\quad Z_0(0)\ne0,\quad
 \int_{\R^2}F'(U)Z_0\dd=2\pi.
\end{equation}
There is a positive regular solution of
\begin{equation}\label{cs:eq:phi}
 \Phi''+r^{-1}\Phi'-4r^{-2}\Phi+F'(U)\Phi=0
\end{equation}
normalized by
\[
 \Phi(r)=r^2+\mathfrak b r^{-2}+O(r^{-4})\quad(r\to\infty).
\]
Its coefficient is
\begin{equation}\label{cs:eq:beta}
 \mathfrak b=\frac14\int_0^\infty r^3F'(U)\Phi\,dr>0.
\end{equation}
\end{lemma}

\begin{proof}
Differentiating the mass identity and the radial equation gives
\eqref{cs:eq:radial-kernel}. In particular the radial kernel is
unbounded, so Lemma~\ref{thm:linearized-classification} gives the
bounded-kernel assertion.

For the regular second mode, \eqref{eq:higher-mode-comparison} gives
$\Phi/\psi>0$ and $(\Phi/\psi)'>0$, where $\psi=-U'>0$.
Its growing $r^2$ coefficient is positive: otherwise $\Phi$ changes
sign or $\Phi/\psi$ tends to zero. Normalize this coefficient to one.
The decay of $F'(U)$ gives the displayed expansion for $\Phi$.
Integrating
\((r^3\Phi'-2r^2\Phi)'=-r^3F'(U)\Phi\)
gives the formula for $\mathfrak b$. Moreover,
$d F(U)/dr=-F'(U)\psi$, and therefore
\begin{align*}
 4\mathfrak b
 &=-\left[r^3F(U)\frac{\Phi}{\psi}\right]_0^\infty
   +\int_0^\infty F(U)
              \left(r^3\frac{\Phi}{\psi}\right)'dr =\int_0^\infty F(U)
       \left\{3r^2\frac{\Phi}{\psi}
                +r^3\left(\frac{\Phi}{\psi}\right)'\right\}dr>0.
\end{align*}
Here $F(U)>0$, $\Phi/\psi>0$, and $(\Phi/\psi)'>0$.
The boundary term vanishes: it is $O(r^4)$ at zero
($\Phi=O(r^2)$, $\psi\asymp r$) and $O(r^{-2})$ at infinity
($F(U)=O(r^{-8})$, $\Phi\asymp r^2$, $\psi\asymp r^{-1}$).
\end{proof}

\subsection{The approximate solution and its residual}

Fix a compact interval $\mathcal D\Subset(0,\infty)$ containing $d_*$ in its
interior. Throughout the construction,
\begin{equation}\label{cs:eq:parameters}
 \delta=d\eps^{2/3},\quad d\in\mathcal D,\quad
 \eta=\eps/\delta,\quad L=1+|\log\eps|,\quad q_\pm=\pm\delta e.
\end{equation}
Here $d$ denotes scalar separation, whereas $\mathbf d$ in
Section~\ref{sec:pointwise} denotes a center shift.
Thus $\delta\asymp\eta^2$; all estimates below are uniform for $d\in\mathcal D$.

Choose a fixed $0<r_c<1/100$ and a smooth radial cutoff $\chi$ with
$0\le\chi\le1$, $\chi=1$ on $[0,1]$, and $\chi=0$ on $[2,\infty)$.
Set $\chi_j(x)=\chi(|x-q_j|/(r_c\delta))$. For opposite indices $j,l$ put
\begin{align}
 B_j(x)&=u_0(x)+16\pi R(x-q_j)+16\pi G(x,q_l),\label{cs:eq:B}\\
 C_{\eps,\delta}&=8\log\eps+c_8-B_j(q_j),\label{cs:eq:C}\\
 O(x)&=C_{\eps,\delta}+u_0(x)+16\pi G(x,q_+)+16\pi G(x,q_-),\label{cs:eq:outer}\\
 I_j(x)&=U((x-q_j)/\eps)+B_j(x)-B_j(q_j),\label{cs:eq:inner}\\
 u^{\mathrm{app}}_{\eps,\delta}&=O+\sum_{j=\pm}\chi_j(I_j-O).\label{cs:eq:app}
\end{align}
Reflection makes \eqref{cs:eq:C} independent of $j$.
The poles at $q_j$ cancel, so $u^{\mathrm{app}}_{\eps,\delta}-u_0$ is smooth and periodic.

Since $\Delta B_j=0$ near $q_j$, the inner residual contains only the
nonlinear change caused by $B_j-B_j(q_j)$. Writing $h_0=u_0-4\log|x|$
near zero, set
\(\calA=D^2h_0(0)+32\pi D^2R(0) =24\pi\{D^2R(0)-D^2G(a,0)\}.\)
The singular linear terms cancel:
\begin{equation}\label{cs:eq:Bgradient}
 \nabla B_+(q_+)=\delta\calA e+O(\delta^3).
\end{equation}
Consequently, for $r=|z|\le r_c/\eta$,
\begin{equation}\label{cs:eq:Bbound}
 |B_j(q_j+\eps z)-B_j(q_j)|
       \le C(\eps\delta r+\eta^2r^2)
       \le C\eta^2(1+r^2).
\end{equation}
Also $C_{\eps,\delta}=8\log\eps+4\log\delta+O(1)$.

Use the weight and norm
\begin{equation}\label{cs:eq:weight}
 \omega_\eps(x)=\sum_{j=\pm}
       \eps^{-2}(1+\dist(x,q_j)/\eps)^{-5},
 \qquad \|h\|_*=\sup_\T |h|/\omega_\eps.
\end{equation}
In particular $\int_\T\omega_\eps\le C$, and for $s\ge3\delta$,
\begin{equation}\label{cs:eq:weight-tail}
 \int_{\T\setminus B_s(0)}\omega_\eps\le C(\eps/s)^3.
\end{equation}

\begin{lemma}\label{cs:lem:residual}
Let
\(
 E=\Delta u^{\mathrm{app}}_{\eps,\delta}+\eps^{-2}F(u^{\mathrm{app}}_{\eps,\delta})
                  -8\pi\delta_0-24\pi\delta_a.
\) Then
\begin{equation}\label{cs:eq:residual}
 \|E\|_*\le C\eta^2,\qquad
 \|\eps^{-2}e^{u^{\mathrm{app}}_{\eps,\delta}}\|_*\le C.
\end{equation}
Moreover,
\begin{equation}\label{cs:eq:app-tail}
 \int_{\T\setminus\bigcup_jB_{r_c\delta}(q_j)}
                \eps^{-2}e^{u^{\mathrm{app}}_{\eps,\delta}}\dd\le C\eta^6,
 \qquad \sup_\T u^{\mathrm{app}}_{\eps,\delta}\longrightarrow U(0)<0.
\end{equation}
\end{lemma}

\begin{proof}
In $B_{r_c\delta}(q_j)$, \eqref{cs:eq:Bbound} and \eqref{cs:eq:tails} bound
the residual by
$C\eps^{-2}\eta^2(1+r)^{-6}$.
On the matching annulus,
\(I_j-O=U(|x-q_j|/\eps)+8\log|x-q_j|-8\log\eps-c_8.\)
Its derivatives of order $l=0,1,2$ are bounded by
$C\delta^{-l}\eta^6$. The cutoff error is
$O(\delta^{-2}\eta^6)=O(\eps^{-2}\eta^8)$, whereas
$\omega_\eps\ge c\eps^{-2}\eta^5$ there.

Outside the core balls and with $|x|\le3\delta$,
$e^{u^{\mathrm{app}}_{\eps,\delta}}\le C\eta^8$. For $3\delta\le|x|\le r_0$,
\(e^{u^{\mathrm{app}}_{\eps,\delta}}\le C\eps^8\delta^4|x|^{-12},\)
and away from zero the bound is $C\eps^8\delta^4$.
The density vanishes at the extra vortex. These bounds prove
\eqref{cs:eq:residual} and the integral estimate in \eqref{cs:eq:app-tail}.
Inner convergence to $U$ and the same tail bounds give the supremum limit.
\end{proof}

\subsection{A uniform projected inverse for the collapsing pair}\label{cs:sec:inverse}

Let $\mathcal X_{\rm sym}$ be the space of functions even in both torus
coordinates. Fix $R_0>1$, independently of $\eps$, and define
\begin{equation}\label{cs:eq:K}
 Z=\partial_eU,\qquad
 K_j(x)=\eps^{-2}\chi(|x-q_j|/(R_0\eps))Z((x-q_j)/\eps),
 \qquad K=K_+-K_- .
\end{equation}
Then $K\in\mathcal X_{\rm sym}$ and $\int K_j=0$ separately.
Write
\(\mathcal J=\int_{\R^2}\chi(|z|/R_0)Z(z)^2\,dz>0.\)
The orthogonality condition is $\int_\T K\phi=0$.
For an even function it removes the surviving vertical translation
at either core; horizontal translations are odd in the first coordinate.

\begin{lemma}\label{cs:lem:inverse}
For all small $\eps$, the problem
\begin{equation}\label{cs:eq:linear}
 [\Delta+\eps^{-2}F'(u^{\mathrm{app}}_{\eps,\delta})]\phi=h+tK,\qquad
 \phi\in\mathcal X_{\rm sym},\qquad \int_\T K\phi=0
\end{equation}
has a unique pair $(\phi,t)$ for every smooth $h\in\mathcal X_{\rm sym}$.
Uniformly for $d\in\mathcal D$,
\begin{equation}\label{cs:eq:inverse}
 \|\phi\|_\infty\le CL\|h\|_*,\qquad |t|\le C\|h\|_*.
\end{equation}
For continuous $h$, the same estimates hold for the unique
distributional solution $\phi\in W^{2,p}(\T)$ for every
$1<p<\infty$, by approximation and elliptic regularity.
\end{lemma}

\begin{proof}
Adapting \cite[Appendix~A]{LinYanAIHP}, we explicitly control
radial averages, the global integral, and all intermediate scales.
Consider the radial reference potential
\(
 V_\eps^0=\sum_j\chi_j(x)\eps^{-2}F'(U((x-q_j)/\eps))
\)
and the corresponding projected problem
\begin{equation}\label{cs:eq:reference-linear}
 (\Delta+V_\eps^0)\phi=h+tK,\qquad
 \phi\in\mathcal X_{\rm sym},\qquad \int_\T K\phi=0.
\end{equation}
Until the final perturbation step, $(\phi,t)$ denotes a solution of
\eqref{cs:eq:reference-linear}.
Testing this equation against $\chi_jZ((x-q_j)/\eps)$ gives
\begin{equation}\label{cs:eq:t-linear}
 |t|\le C\{\|h\|_*+\eta\|\phi\|_\infty\}.
\end{equation}
Indeed, $2R_0\eps<r_c\delta$ for small $\eps$, so the test
function equals $Z$ on the support of $K_j$ and the coefficient of
$t$ has absolute value $\mathcal J$. Writing $Z_j(x)=Z((x-q_j)/\eps)$,
the reference operator applied to the test function is
\[
 (\Delta\chi_j)Z_j+2\nabla\chi_j\cdot\nabla Z_j
       +\eps^{-2}(\chi_j^2-\chi_j)F'(U((x-q_j)/\eps))Z_j.
\]
All terms are supported where $|x-q_j|\asymp\delta$.
The first two have $L^1$ norm $O(\eta)$ since
$Z_j=O(\eta)$ and $\nabla Z_j=O(\eps/\delta^2)$; the last has
$L^1$ norm $O(\eps^{-2}\delta^2\eta^9)=O(\eta^7)$ by
\eqref{cs:eq:tails}. This proves \eqref{cs:eq:t-linear}.

We prove the sup-norm estimate by contradiction. Otherwise normalize
$\|\phi\|_\infty=1$ and let $L\|h\|_*\to0$. By
\eqref{cs:eq:t-linear}, $t\to0$. The inner limits solve the bounded
kernel equation. Lemma~\ref{cs:lem:profile}, parity, and orthogonality
show that these limits are zero.

We give the quantitative radial estimate before treating the
intermediate scales. Write $H_* = \|h\|_*$,
$P=\|\phi\|_\infty$, and $R_\eta=r_c/\eta$. For $0\le s\le R_\eta$, set
\(
 y_j(s)=\frac1{2\pi}\int_0^{2\pi}
       \phi(q_j+\eps s(\cos\vartheta,\sin\vartheta))\,d\vartheta.
\)
On this disk the cutoff is one and the other core's potential and
projection are absent. The projection at $q_j$ has zero circular
average. Hence, with $f_j$ the circular average of $\eps^2h$,
\begin{equation}\label{cs:eq:radial-forcing}
 y_j''+s^{-1}y_j'+F'(U)y_j=f_j(s),\qquad
 |f_j(s)|\le CH_*(1+s)^{-5}.
\end{equation}
The contribution from the other center in $\omega_\eps$ obeys the
same bound because $s\le r_c/\eta$.

Choose a fundamental solution $Y_0$ with
$Z_0Y_0'-Z_0'Y_0=1/s$, for example by setting
$Y_0(s_*)=0$ and $Y_0'(s_*)=1/(s_*Z_0(s_*))$ at any
$s_*>0$ where $Z_0(s_*)\ne0$. Continue it by the regular ODE
across the simple zeros of $Z_0$. This avoids integrating the
reduction-of-order expression through a pole.
Near zero,
$Y_0(s)=Z_0(0)^{-1}\log s+O(1)$.
At infinity, $F'(U)=O(s^{-8})$ implies that both solutions are
linear combinations of $1+o(1)$ and $\log s+O(1)$.

With $a_{0,j}=y_j(0)/Z_0(0)$, the globally valid Volterra formula is
\begin{equation}\label{cs:eq:radial-volterra}
 \begin{split}
 y_j(s)&=a_{0,j}Z_0(s)+\mathcal T_j(s),\\
 \mathcal T_j(s)&=Y_0(s)\int_0^s tZ_0(t)f_j(t)\,dt
                  -Z_0(s)\int_0^s tY_0(t)f_j(t)\,dt.
 \end{split}
\end{equation}
Both integrals converge at zero, and this formula contains no
division by $Z_0(s)$. For $0<t\le s$, its kernel satisfies
\[
 |Y_0(s)Z_0(t)-Z_0(s)Y_0(t)|
   \le C\log(2+s)\{1+|\log t|+\log(2+t)\}.
\]
Integrating against $t(1+t)^{-5}$ proves
\begin{equation}\label{cs:eq:radialaverage}
 \bar\phi_j(r)=a_{0,j}Z_0(r/\eps)+T_j(r),\qquad
 |T_j(r)|\le CH_*\log(2+r/\eps),\quad r\le r_c\delta,
\end{equation}
where $T_j(r)=\mathcal T_j(r/\eps)$. Moreover,
\begin{equation}\label{cs:eq:radial-weighted-error}
 \int_0^{R_\eta}|F'(U(s))\mathcal T_j(s)|s\,ds\le CH_*.
\end{equation}
This bound is uniform in the growing upper limit, since
$\int_0^\infty s(1+s)^{-8}\log(2+s)\,ds<\infty$.

Reflection interchanges the cores, so $a_{0,+}=a_{0,-}=a_0$.
Integrate the reference equation on the torus. Periodicity and
$\int K=0$ give exactly $\int V_\eps^0\phi=\int h$.
Separating the disks where $\chi_j=1$ from the cutoff annuli gives
\begin{equation}\label{cs:eq:radial-torus-identity}
 \int_\T h
   =4\pi a_0\int_0^{R_\eta}F'(U(s))Z_0(s)s\,ds
     +2\pi\sum_{j=\pm}\int_0^{R_\eta}
                     F'(U(s))\mathcal T_j(s)s\,ds+E_{\rm cut},
\end{equation}
where
\[
 |E_{\rm cut}|\le CP\int_{R_\eta}^{2R_\eta}s(1+s)^{-8}\,ds
       \le C\eta^6P.
\]
Outside the cutoff disks the reference potential is exactly zero.
Also $|\int h|\le CH_*$, and \eqref{cs:eq:radial-kernel} gives
\[
 \int_0^{R_\eta}F'(U)Z_0s\,ds=1+O(\eta^6L).
\]
Thus \eqref{cs:eq:radial-torus-identity} yields
\[
 4\pi a_0=O(H_*+\eta^6P+|a_0|\eta^6L).
\]
Absorb the last term, using $\eta^6L\to0$, to obtain
\begin{equation}\label{cs:eq:radial-quantitative}
 |a_0|\le C(H_*+\eta^6P),\qquad
 \sup_{r\le r_c\delta}|\bar\phi_j(r)|
       \le CL(H_*+\eta^6P).
\end{equation}
In the contradiction sequence $P=1$ and $LH_*\to0$, so
\begin{equation}\label{cs:eq:avgozero}
 \sup_{r\le r_c\delta}|\bar\phi_j(r)|\longrightarrow0.
\end{equation}
Here reflection makes the two radial coefficients equal, which is
what permits the global identity to control each one.

On the $\delta$ scale the rescaled reference potential is
$O(\eta^6)$ and the rescaled forcing is $O(H_*\eta^3)$ on compact
sets away from $\pm e$; the projection supports shrink to these
two points. Every limit is therefore bounded and harmonic on
$\R^2\setminus\{\pm e\}$. Removability and Liouville's theorem
make it constant. Its average on a circle of fixed radius less
than $r_c$ around either core is zero by \eqref{cs:eq:avgozero}.
Thus the cluster-scale limit is zero.

The outer logarithmic mode is ruled out by an exact flux identity.
For $3\delta\le r\le r_0$, all reference potentials and projection
terms lie in $B_r(0)$. If $\bar\phi_0$ denotes the circular average
about zero, subtraction of the integrated torus equation gives
\begin{equation}\label{cs:eq:outer-flux}
 2\pi r\bar\phi_0'(r)
      =\int_{B_r(0)}h-\int_\T h
      =-\int_{\T\setminus B_r(0)}h.
\end{equation}
Using \eqref{cs:eq:weight-tail} and integrating in $r$ yields
\begin{equation}\label{cs:eq:outer-average}
 |\bar\phi_0(r)-\bar\phi_0(3\delta)|
     \le CH_*\eps^3\int_{3\delta}^{r}s^{-4}\,ds
     \le CH_*\eta^3.
\end{equation}
The average at $3\delta$ tends to zero by the cluster-scale limit.
Consequently all these outer averages tend uniformly to zero;
an independent term proportional to $\log r$ is excluded by
\eqref{cs:eq:outer-flux}.

To pass from local limits to the global supremum, choose points
$x_\eps$ with $|\phi(x_\eps)|=1$. After taking a subsequence,
exactly one of the following alternatives applies:
\begin{enumerate}[label=(\roman*),leftmargin=*]
\item $\dist(x_\eps,\{q_+,q_-\})=O(\eps)$.
The inner limit is zero by the bounded-kernel classification,
parity, and orthogonality proved above.
\item $s=|x_\eps-q_j|$ satisfies $\eps\ll s\ll\delta$ for one core.
After scaling about $q_j$ by $s$, the potential on compact annuli
is $O((\eps/s)^6)$ and the forcing is $O(H_*(\eps/s)^3)$.
The other core leaves every compact set and the projection
contracts to the origin. A bounded harmonic limit on
$\R^2\setminus\{0\}$ is constant, and its unit-circle average is
zero by \eqref{cs:eq:avgozero}.
\item $|x_\eps|=O(\delta)$ and its distance from both cores is
bounded below by a positive multiple of $\delta$.
The cluster-scale limit just established is zero.
\item $s=|x_\eps|$ satisfies $\delta\ll s\ll1$.
The reference potential and projection are supported in a disk
of radius $O(\delta/s)$ after scaling. On compact annuli the
forcing is $O(H_*(\eps/s)^3)$. The bounded harmonic limit on
$\R^2\setminus\{0\}$ is constant with zero average by
\eqref{cs:eq:outer-average}.
\item $|x_\eps|$ stays bounded away from zero.
The limiting function is bounded and harmonic on the punctured
torus, extends across zero, and is constant. Its average on any
fixed circle of radius less than $r_0$ is zero by
\eqref{cs:eq:outer-average}.
\end{enumerate}
In each case elliptic compactness holds near the rescaled maximum
point, which lies away from the listed punctures, so its value
converges to zero. This contradicts $|\phi(x_\eps)|=1$ and proves
$\|\phi\|_\infty\le CLH_*$ for the reference operator.
Equation~\eqref{cs:eq:t-linear} and $\eta L\to0$ give $|t|\le CH_*$.

The augmented map
$(\phi,t)\mapsto((\Delta+V_\eps^0)\phi-tK,\int K\phi)$
is Fredholm of index zero on the even periodic $H^2$ and $L^2$
spaces with one added scalar variable. The estimate makes its
kernel trivial; it is therefore invertible.

We now return to the potential in \eqref{cs:eq:linear}.
The same inner and annular bounds used in Lemma~\ref{cs:lem:residual} give
\(
 \|(\eps^{-2}F'(u^{\mathrm{app}}_{\eps,\delta})-V_\eps^0)\phi\|_*
                      \le C\eta^2\|\phi\|_\infty .
\)
Since $\eta^2L\to0$, a Neumann-series perturbation of the reference
augmented inverse proves \eqref{cs:eq:linear}--\eqref{cs:eq:inverse}.
\end{proof}

\subsection{The projected nonlinear solution and its mass}

\begin{proposition}\label{cs:prop:projected}
There is $C_0>0$ such that, for every $d\in\mathcal D$ and small
$\eps$, there is a unique correction $\phi=\phi_{\eps,d}\in\mathcal X_{\rm sym}$
in the ball $\|\phi\|_\infty\le C_0\eta^2L$, with
$\int K\phi=0$, for which $u=u^{\mathrm{app}}_{\eps,\delta}+\phi$ solves
\begin{equation}\label{cs:eq:projected}
 \Delta u+\eps^{-2}F(u)
       =8\pi\delta_0+24\pi\delta_a+t_{\eps,d}K.
\end{equation}
It depends continuously on $d$ and satisfies
\begin{equation}\label{cs:eq:phi-bound}
 \|\phi\|_\infty\le C\eta^2L,\qquad |t_{\eps,d}|\le C\eta^2.
\end{equation}
There is a fixed $\gamma>0$ such that $u\le-\gamma$, and
\begin{equation}\label{cs:eq:mass}
 \int_\T\eps^{-2}F(u)\dd=32\pi,\qquad
 \int_{\T\setminus\bigcup_jB_{r_c\delta}(q_j)}
                     \eps^{-2}F(u)\dd=O(\eta^6).
\end{equation}
\end{proposition}

\begin{proof}
Write $u^{\mathrm{app}}=u^{\mathrm{app}}_{\eps,\delta}$. The equation for $\phi$ has right-hand side
$-E-N(\phi)+tK$, where
\[
 N(\phi)=\eps^{-2}\{F(u^{\mathrm{app}}+\phi)-F(u^{\mathrm{app}})-F'(u^{\mathrm{app}})\phi\}.
\]
For $\|\phi\|_\infty,\|\widetilde\phi\|_\infty\le o(1)$,
\[
 \|N(\phi)\|_*\le C\|\phi\|_\infty^2,\quad
 \|N(\phi)-N(\widetilde\phi)\|_*
 \le C(\|\phi\|_\infty+\|\widetilde\phi\|_\infty)
                                      \|\phi-\widetilde\phi\|_\infty .
\]
For large fixed $C_0$, Lemmas~\ref{cs:lem:residual} and
\ref{cs:lem:inverse} give a self-map contraction with constant
$O(\eta^2L^2)=o(1)$ and the bound for $t$.
Uniqueness and continuity of the coefficients and augmented inverse
give continuity in $d$. Elliptic regularity gives $u-u_0\in C^\infty$,
while Lemma~\ref{cs:lem:residual} gives $u\le-\gamma$.
Integrating \eqref{cs:eq:projected} and using $\int K=0$ gives the exact
mass; \eqref{cs:eq:app-tail} and $\|\phi\|_\infty=o(1)$ give the tail bound.
\end{proof}

\subsection{Actual masses, centers, and quadrupole moments}\label{cs:sec:moments}

In this section $u$ denotes the projected solution.
In the coordinate whose real axis is $e$, define the effective measure
\begin{equation}\label{cs:eq:nu}
 d\nu=\frac1{2\pi}\{\eps^{-2}F(u)-tK\}\dd,\quad
 D_j=B_{r_c\delta}(q_j),\quad \nu_j=\nu|_{D_j},\quad
 \nu_{\rm out}=\nu-\nu_+-\nu_-.
\end{equation}
The outside measure is positive; the signed core measures have
uniformly bounded total variations. Proposition~\ref{cs:prop:projected}
and $\int_{D_j}K_j=0$ give their positive masses:
\begin{equation}\label{cs:eq:actual-mass}
 \nu(\T)=16,\qquad
 \widetilde m:=\nu(D_+)=\nu(D_-)=8+O(\eta^6),\qquad
 \nu_{\rm out}(\T)=O(\eta^6).
\end{equation}

All complex moments below are real because of reflection in the
transverse coordinate. The original-center first moments satisfy
\begin{equation}\label{cs:eq:firstmoment}
 \left|\int_{D_j}(x-q_j)\,d\nu_j(x)\right|\le C\eps\eta^2L.
\end{equation}
Indeed, the radial measure
\(d\nu_j^0=(2\pi)^{-1}\eps^{-2} F(U((x-q_j)/\eps))1_{D_j}\dd\)
has zero nonconstant complex moments, and
\begin{equation}\label{cs:eq:density-difference}
 |F(u(q_j+\eps z))-F(U(z))|
       \le C(1+|z|)^{-8}\eta^2(|z|^2+L).
\end{equation}
The first moment of the projection has size $O(\eps|t|)$.
This proves \eqref{cs:eq:firstmoment}.
The same density bound also gives
\begin{equation}\label{cs:eq:second-absolute-moment}
 \int_{D_j}|x-q_j|^2\,d|\nu_j|(x)\le C\eps^2.
\end{equation}
Indeed, the radial term has a finite second absolute moment,
the density difference contributes
$C\eps^2\eta^2\int_0^{r_c/\eta}r^3(1+r)^{-8}(r^2+L)\,dr
\le C\eps^2\eta^2L$, and the projection contributes
$O(\eps^2|t|)$. The estimate remains valid about $\zeta_j$ below.

Recenter each signed core measure at its center of mass:
\begin{equation}\label{cs:eq:centers}
 \zeta_j=q_j+\frac1{\widetilde m}\int_{D_j}(x-q_j)d\nu_j(x),\qquad
 \zeta_\pm=\pm\widehat\delta e,\quad
 |\widehat\delta-\delta|\le C\eps\eta^2L=o(\eps).
\end{equation}
These expansion centers are well defined since $\widetilde m$ stays positive. Let
\(M_{j,n}=\int_{D_j}(x-\zeta_j)^n\,d\nu_j(x).\)
Then $M_{j,1}=0$ and $M_{-,n}=(-1)^nM_{+,n}$.

\begin{lemma}\label{cs:lem:moments}
Uniformly for $d\in\mathcal D$,
\begin{align}
 M_{j,2}&=-2\mathfrak b\frac{\eps^4}{\delta^2}
      +O\left(\eps^2\eta^4L^2+\eps^4\right),\label{cs:eq:secondmoment}\\
 \sum_{n\ge3}\frac{|M_{j,n}|}{\widehat\delta^{\,n+1}}
                       &\le C\frac{\eta^5L}{\delta}.\label{cs:eq:highermom}
\end{align}
\end{lemma}

\begin{proof}
Write $\upsilon(z)=u(q_++\eps z)-U(z)$, $R_\eta=r_c/\eta$, and
\[
 \upsilon_2(r)=\frac1\pi\int_0^{2\pi}\upsilon(r,\theta)\cos2\theta\,d\theta.
\]
On $r\le R_\eta$, \eqref{cs:eq:Bbound} and \eqref{cs:eq:phi-bound} give
\(|\upsilon|\le C\eta^2(r^2+L).\)
Let $n_2$ be the mode-two coefficient of
$F(U+\upsilon)-F(U)-F'(U)\upsilon$. The projected forcing has only angular mode one
at this core, so
\begin{equation}\label{cs:eq:mode2}
 \upsilon_2''+r^{-1}\upsilon_2'-4r^{-2}\upsilon_2+F'(U)\upsilon_2=-n_2 .
\end{equation}
Taylor's theorem gives
\begin{equation}\label{cs:eq:n2}
 |n_2(r)|\le C\eta^4(r^{-4}+L^2r^{-8})\quad(r\ge1).
\end{equation}
For $r\le1$ the bound is $C\eta^4L^2r^2$. To justify the factor
$r^2$, interior elliptic estimates in the projected equation give
$\|\upsilon\|_{C^2(B_2)}\le C\eta^2L$; the nonlinear Taylor remainder has
$C^2$ norm $O(\eta^4L^2)$, and its second angular mode vanishes to
order two at the origin.

The harmonic incident field $B_+$ has quadratic coefficient
\(a_2^{\,0}=-\delta^{-2}+O(1).\)
Since the cutoff in \eqref{cs:eq:app} is one on $\partial D_+$,
\(\upsilon_2(R_\eta)=\eps^2a_2^{\,0}R_\eta^2+O(\eta^2L).\)
Write $\upsilon_2=A_\eta\Phi+p_\eta$, where $p_\eta$ is regular at zero and
vanishes at $R_\eta$. It follows that
\begin{equation}\label{cs:eq:Aeta}
 A_\eta=-\eta^2+O(\eta^4L+\eta^2\delta^2).
\end{equation}
Here $\Phi(R_\eta)/R_\eta^2=1+O(\eta^4)$.

For completeness, variation of parameters for \eqref{cs:eq:mode2} can
be written using the positive solution $\Phi$:
\[
 p_\eta(r)=
 \Phi(r)\int_r^{R_\eta}\frac1{s\Phi(s)^2}
                  \int_0^s t\Phi(t)n_2(t)\,dt\,ds .
\]
The estimates $\Phi(r)\asymp r^2$ at both endpoints and
\eqref{cs:eq:n2} imply
\[
 |p_\eta(r)|\le
 \begin{cases}
 C\eta^4L^2,&0\le r\le1,\\
 C\eta^4r^{-2}\{L^2+\log(2+r)\},&1\le r\le R_\eta.
 \end{cases}
\]
Consequently
\[
 \int_0^{R_\eta}r^3
       \bigl(|F'(U)p_\eta|+|n_2|\bigr)\,dr\le C\eta^4L^2.
\]
Using \eqref{cs:eq:beta}, \eqref{cs:eq:Aeta}, and
$\int_{R_\eta}^\infty r^3F'(U)\Phi\,dr=O(\eta^2)$ gives
\[
 \int_{D_+}(x-q_+)^2\,d\nu_+(x)
 =\frac{\eps^2}{2}\int_0^{R_\eta}r^3
                      \{F'(U)\upsilon_2+n_2\}\,dr
 =-2\mathfrak b\eps^2\eta^2
        +O\bigl(\eps^2\eta^4L^2+\eps^2\eta^2\delta^2\bigr).
\]
The projection has zero second complex moment about $q_+$.
Recentering changes this expression by
$O(\eps^2\eta^4L^2)$, by \eqref{cs:eq:firstmoment}. This proves
\eqref{cs:eq:secondmoment}; reflection treats the other core.

For the higher moments, \eqref{cs:eq:density-difference} and
$|t|\le C\eta^2$ imply
\begin{equation}\label{cs:eq:thirdabsolute}
 \int_{D_j}|x-q_j|^3\,d|\nu_j-\nu_j^0|
                       \le C\eps^3\eta^2L.
\end{equation}
The integral in the scaled variable is finite since
$\int_1^\infty r^6(1+r)^{-8}dr<\infty$.
The radial measure $\nu_j^0$ has no nonconstant complex moments
about $q_j$; about $\zeta_j$ its $n$th moment is its mass times
$(q_j-\zeta_j)^n$. All supports lie in a disk of radius
$2r_c\delta$ about $\zeta_j$. Summing the geometric series beginning
at order three, using \eqref{cs:eq:thirdabsolute} and
\eqref{cs:eq:centers}, therefore gives
\[
 \sum_{n\ge3}\frac{|M_{j,n}|}{\widehat\delta^{\,n+1}}
 \le C\widehat\delta^{-4}
       \{\eps^3\eta^2L+|q_j-\zeta_j|^3\}
 \le C\frac{\eta^5L}{\delta}.
\]
\end{proof}

\subsection{The translation equation with a controlled error}

Green representation for \eqref{cs:eq:projected} is exact:
\begin{equation}\label{cs:eq:representation}
 u(x)=u_0(x)+2\pi\int_\T G(x,y)\,d\nu(y)+\overline c .
\end{equation}
Consider the positive core and the circle
$|x-\zeta_+|=4r_c\delta$. The core supports lie inside radius
$2r_c\delta$ about their respective centers.
On this circle the outside-core density and its potential obey
\begin{equation}\label{cs:eq:neck-potential}
 \eps^{-2}|\mathcal P(u)|\le C\delta^{-2}\eta^6,\qquad
 \left|\nabla\left(2\pi\int G(x,y)d\nu_{\rm out}(y)\right)\right|
                         \le C\eta^6/\delta.
\end{equation}
For the second estimate, integrate the locally integrable logarithmic
gradient against the cluster density $C\delta^{-2}\eta^6$ and the outer
bound $C\eps^6\delta^4|y|^{-12}$ from Lemma~\ref{cs:lem:residual}.
Each contributes $O(\eta^6/\delta)$; the distant part is smaller.

After subtracting this potential, the remaining field has constant
Laplacian $4\pi(\widetilde m-8)$ on the matching annulus.
In a complex coordinate $w$ centered at $\zeta_+$ and directed along
$e$, subtracting $\pi(\widetilde m-8)|w|^2$ therefore gives a harmonic
field with expansion
\begin{equation}\label{cs:eq:Laurent}
 H(w)=-\widetilde m\log|w|
             +\operatorname{Re}\sum_{n\ge1}(a_nw^n+b_nw^{-n}),
 \qquad b_n=M_{+,n}/n .
\end{equation}
An additive constant has been suppressed.
The quadratic correction accounts for the regular Green parts when
only $\nu_++\nu_-$ is retained. Its gradient on the matching circle is
$O(\eta^6\delta)$ and is absorbed in \eqref{cs:eq:neck-potential}.
Thus $|\nabla u-\nabla H|\le C\eta^6/\delta$ there.
All coefficients are real, and $b_1=0$.

We compute the incident field term by term. Put
$D=\widehat\delta$ and $\Delta m=\widetilde m-8=O(\eta^6)$.
For $|w|<D$, the singular incident part, up to a constant, is
\[
 \operatorname{Re}\left\{\begin{aligned}
 &4\log(1+w/D)-8\log(1+w/(2D))\\
 &\quad-\Delta m\log(1+w/(2D))+\frac{b_2}{(2D+w)^2}
       +\sum_{n\ge3}\frac{M_{-,n}}{n(2D+w)^n}
 \end{aligned}\right\}.
\]
The signs of the moments here follow from expanding
$-\int\log|2D+w-(x-\zeta_-)|\,d\nu_-(x)$.
In particular $M_{-,2}=M_{+,2}=2b_2$. The relevant coefficients are
\[
\begin{array}{c|ccc}
 \text{term} & [w] & [w^2] & [w^3]\\ \hline
 4\log(1+w/D) & 4/D & -2/D^2 & 4/(3D^3)\\
 -8\log(1+w/(2D)) & -4/D & 1/D^2 & -1/(3D^3)\\
 -\Delta m\log(1+w/(2D))
      & -\Delta m/(2D)&\Delta m/(8D^2)&-\Delta m/(24D^3)\\
 b_2/(2D+w)^2 &-b_2/(4D^3)&3b_2/(16D^4)&-b_2/(8D^5)
\end{array}
\]
Thus the two leading logarithms cancel in degree one and have
cubic coefficient $D^{-3}$.

Replacing both cores by masses $8$ at $\pm De$ in the regular Green
field gives linear coefficient $\lambda D+O(D^3)$ by
\eqref{cs:eq:Bgradient} and cubic coefficient $O(1)$.
The actual-mass correction is $O(\eta^6)$ in each fixed derivative.
The source derivatives are uniformly bounded; hence zero centered
dipoles and \eqref{cs:eq:second-absolute-moment} give a Taylor error
$O(\eps^2)=O(\delta^3)$, using $\eps^2=\delta^3/d^3$.
Also, $\Delta R=1$ makes $\nabla R$ componentwise harmonic, so radial
measures act exactly as point masses.

For degree one the moments of order $n\ge3$ contribute at most
$C\sum_{n\ge3}|M_{-,n}|/D^{n+1}
\le C\eta^5L/\delta$.
For degree three their coefficients contain
$\binom{n+2}{3}/(n2^{n+3})$, which is bounded uniformly in $n$.
Consequently their total cubic contribution is
$O(\eta^5L/\delta^3)$.
The quadrupole has $|b_2|/D^2=O(\eta^4)$ by
\eqref{cs:eq:secondmoment}, while the mass correction contributes
$O(\eta^6/\delta)$ and $O(\eta^6/\delta^3)$ to the linear and
cubic coefficients, respectively. Combining these separate terms,
\begin{align}
 a_1={}&\lambda D-\frac{b_2}{4D^3}
       +O\left(\delta^3+\frac{\eta^6}{\delta}
                              +\frac{\eta^5L}{\delta}\right),\label{cs:eq:a1}\\
 a_3={}&D^{-3}\{1+O(\eta^4+\eta^5L+\eta^6+\delta^3)\},\qquad
 |a_n|\le \frac{C}{nD^n}\quad(n\ge1).\label{cs:eq:an}
\end{align}
For the last bound, the vortex is at distance $D$ and the opposite
core at distances at least $(2-3r_c)D>D$. Their logarithmic expansions
and bounded total variations give $C/(nD^n)$ uniformly in $n$.
The smooth field is analytic on a fixed-radius disk, so its
geometric coefficient bounds satisfy the same estimate for small $\eps$.

\begin{lemma}\label{cs:lem:reduced}
With $\mathcal J$ as in Section~\ref{cs:sec:inverse}, uniformly for
$d\in\mathcal D$ as $\eps\to0$,
\begin{equation}\label{cs:eq:reduced}
 \frac{t_{\eps,d}}{\eps^{5/3}}
   =-\frac{16\pi}{\mathcal J}
        \left(\lambda d-\frac{\mathfrak b}{2d^5}
                            +O(\eps^{1/3}L)\right).
\end{equation}
\end{lemma}

\begin{proof}
Use the stress tensor
\[
 \mathsf T(u)=\nabla u\otimes\nabla u-\tfrac12|\nabla u|^2I_2
                              +\eps^{-2}\mathcal P(u)I_2 .
\]
On the positive-core disk,
$\operatorname{div}\mathsf T(u)=tK_+\nabla u$.
With $\boldsymbol n$ the outward unit normal, its flux in the $e$
direction is therefore
\begin{equation}\label{cs:eq:stressrhs}
 \int_{\partial B_{4r_c\delta}(\zeta_+)}\mathsf T(u)\boldsymbol n\cdot e\,ds
       =\frac{t}{\eps}\{\mathcal J+O(\eta^2L)\}.
\end{equation}
Indeed, on the fixed support of the scaled projection,
\(\|u(q_++\eps\,\cdot)-U\|_{C^1(B_{2R_0})}\le C\eta^2L.\)
This follows from \eqref{cs:eq:Bbound}, \eqref{cs:eq:phi-bound},
the bound on $t$, and interior elliptic estimates in the projected
equation. In particular the factor multiplying $t/\eps$ stays
bounded away from zero, uniformly in $d$.

The pressure term on the left is $O(\eta^6/\delta)$ by
\eqref{cs:eq:neck-potential}. Replacing the gradient by that of $H$
has the same error, since $|\nabla H|\le C/\delta$ on the circle.
Write $w=\xi+i\upsilon$ with the $\xi$-axis directed along $e$,
and use the convention $H_w=(H_\xi-iH_\upsilon)/2$.
For the harmonic expansion \eqref{cs:eq:Laurent}, direct Laurent
multiplication gives
\[
 \Res_{w=0}(H_w)^2
   =-\frac12\left(\widetilde m a_1+
                    \sum_{n\ge1}n(n+1)b_na_{n+1}\right).
\]
The stress flux in the $e$ direction is $4\pi$ times this real residue.
The $n=1$ term vanishes. Moreover, \eqref{cs:eq:an} and
\eqref{cs:eq:highermom} give
\(\sum_{n\ge3}n(n+1)|b_na_{n+1}| \le C\eta^5L/\delta=o(\delta).\)
We used $\delta\asymp\eta^2$.
Equations \eqref{cs:eq:actual-mass}, \eqref{cs:eq:secondmoment},
\eqref{cs:eq:a1}, and \eqref{cs:eq:an} now imply
\begin{align*}
 a_1+\frac6{\widetilde m} b_2a_3
  &=\lambda\widehat\delta
       +\left(-\frac14+\frac68\right)\frac{b_2}{\widehat\delta^3}
       +o(\delta)\\
  &=\lambda\widehat\delta
       -\frac{\mathfrak b}{2}\frac{\eps^4}{\widehat\delta^5}+o(\delta).
\end{align*}
We detail uniformity of the discarded terms. The identities
$\delta=d^3\eta^2$ and
$|D/\delta-1|\le C\eta^3L$ hold uniformly on $\mathcal D$.
The force errors relative to the leading size $\delta$ are
\[
\begin{array}{l|c|c}
 \text{source} & \text{force error} & \text{error divided by }\delta\\ \hline
 \text{smooth field}&O(\delta^3)&O(\eps^{4/3})\\
 \text{mass, outer potential, pressure}
     &O(\eta^6/\delta)&O(\eps^{2/3})\\
 \text{moments of order }n\ge3
     &O(\eta^5L/\delta)&O(\eps^{1/3}L)\\
 \text{second-moment remainder}
     &O((\eps^2\eta^4L^2+\eps^4)/\delta^3)
     &O(\eps^{2/3}L^2+\eps^{4/3})\\
 \text{change from }D\text{ to }\delta
     &O(\delta\eta^3L)&O(\eps L)
\end{array}
\]
The second-moment row uses the quantitative remainder in
\eqref{cs:eq:secondmoment}. The relative error in $a_3$ contributes
$O(\eta^4+\eta^5L+\eta^6+\delta^3)$, because
$b_2/D^3=O(\delta)$. Replacing $\widetilde m$ by $8$ contributes
$O(\eta^6)$ relatively, and division by the coefficient in
\eqref{cs:eq:stressrhs} contributes $O(\eta^2L)$ relatively.
Every entry tends to zero and is bounded by $C\eps^{1/3}L$ for
small $\eps$, uniformly in $d\in\mathcal D$.
For recentering, the derivative of
$\lambda D-\mathfrak b\eps^4/(2D^5)$ is bounded when
$D\asymp\delta=d\eps^{2/3}$, which justifies the indicated error.

Combining the stress identity and these estimates therefore gives
\[
 \frac{t}{\eps}
   =-\frac{16\pi}{\mathcal J}
        \left(\lambda\delta-\frac{\mathfrak b\eps^4}{2\delta^5}
                         +O(\delta\eps^{1/3}L)\right).
\]
Substitution of $\delta=d\eps^{2/3}$ proves \eqref{cs:eq:reduced}.
Only uniform convergence and continuity in $d$ are used to choose
the scale; no derivative estimate for the remainder is needed.
\end{proof}

\subsection{The positive periodic coefficient and the choice of scale}

For the square torus, the quarter-turn symmetry and $\Delta R=1$
give $D^2R(0)=\frac12I$. Summing the Green Fourier series in the
second frequency gives, for $y=0$ away from the source,
\[
 G(x,0)=C_0+\sum_{n\ge1}\frac{\coth(\pi n)}{2\pi n}\cos(2\pi nx).
\]
The part with $\coth(\pi n)$ replaced by $1$ is Abel summed; the
difference has exponentially convergent derivatives.
At $a=(1/2,0)$ this yields
\begin{equation}\label{cs:eq:Gxx}
 G_{xx}(a,0)=\frac\pi2
       -4\pi\sum_{n\ge1}\frac{(-1)^n n}{e^{2\pi n}-1}>\frac\pi2>\frac12.
\end{equation}
The alternating-series test gives the strict inequality.
Reflection and the equation give $G_{xy}(a,0)=0$ and
$G_{yy}(a,0)=1-G_{xx}(a,0)$. Thus
$\calA=\operatorname{diag}(-\lambda,\lambda)$ with $\lambda>0$,
establishing the signs in \eqref{cs:eq:lambda} and \eqref{cs:eq:a1}
for the vertical direction.

\begin{proof}[Proof of Theorem~\ref{cs:thm:existence}]
The function $\mathscr G(d)=\lambda d-\mathfrak b/(2d^5)$ has unique
positive zero $d_*$, with $\mathscr G'(d_*)=6\lambda>0$.
Fix $d_-<d_*<d_+$ in $\mathcal D$. By
Proposition~\ref{cs:prop:projected} and Lemma~\ref{cs:lem:reduced},
$t_{\eps,d}$ is continuous and has opposite endpoint signs for every
small $\eps$. The intermediate value theorem gives
$d_\eps\in(d_-,d_+)$ with $t_{\eps,d_\eps}=0$, hence an exact solution
of \eqref{cs:eq:target}. Uniform convergence in \eqref{cs:eq:reduced}
and uniqueness of the zero imply $d_\eps\to d_*$.

Equations \eqref{cs:eq:inner}, \eqref{cs:eq:Bbound},
\eqref{cs:eq:phi-bound}, and elliptic regularity give inner convergence.
Since $D^2U(0)=-\tfrac12 F(U(0))I<0$, local strict concavity gives
maxima at $q_\pm+o(\eps)$, which reflection places on the vertical axis
in an opposite pair. Their positive distance from zero can replace
$\delta$ in \eqref{cs:eq:scale} and \eqref{cs:eq:profiles}.

The nonnegative density has total mass $32\pi$, exhausted by the
two inner limits of mass $16\pi$. This proves \eqref{cs:eq:measure}
and the no-neck assertion; \eqref{cs:eq:mass} gives the cluster-scale bound.
Away from zero the approximation is the outer field and
$C_{\eps,\delta}-2\log\eps
=6\log\eps+4\log\delta+O(1)\to-\infty$.
This proves the asserted behavior of $w_\eps$.
At a core, $u_0(q_\pm)=4\log\delta+O(1)$, so
\[
 w_\eps(q_{\eps,+})+6\log|q_{\eps,+}|
       =2\log(\delta/\eps)+O(1)\longrightarrow+\infty .
\]
The uniform negative bound and regularity were proved in
Proposition~\ref{cs:prop:projected}. This completes the construction.
\end{proof}

\section{Construction of a non-simple mean-field cluster}\label{sec:mf-construction}

We prove Theorem~\ref{mf:thm:main} on the fixed disk $B_1\subset\C$,
with a strength-one vortex at the origin and the boundary traces
in \eqref{mf:eq:boundary}. The seed uses the branched map $z\mapsto z^2$,
as in \cite[Section~2]{DAprileWeiZhangCVPDE}; the perturbation solves
the full Chern--Simons equation \eqref{mf:eq:CS}.

\subsection{The explicit Liouville family}

For $s>0$, set
\begin{equation}\label{mf:eq:seed}
v_s(z)=\log\frac{32s^4}{\bigl(s^4+|z^2-s|^2\bigr)^2},
\qquad H_s(z)=|z|^2e^{v_s(z)}.
\end{equation}
For $f_s(z)=(z^2-s)/s^2$, the identity
\(\Delta\log(1+|f|^2)=\frac{4|f'|^2}{(1+|f|^2)^2}\)
for holomorphic $f$ gives
\begin{equation}\label{mf:eq:seedPDE}
-\Delta v_s=H_s,
\qquad H_s=\frac{8|f_s'|^2}{(1+|f_s|^2)^2}.
\end{equation}
Since $f_s$ has degree two, the change-of-variables formula with
multiplicity yields
\begin{equation}\label{mf:eq:totalmass}
\int_{\R^2}H_s\dd
=2\int_{\R^2}\frac8{(1+|w|^2)^2}\,dw=16\pi.
\end{equation}

The maxima of $v_s$ are $q_\pm(s)=\pm\sqrt{s}$. Define
\begin{equation}\label{mf:eq:scales}
\delta_s=\sqrt{s},\qquad \ell_s=\frac{s^{3/2}}2.
\end{equation}
Then $\ell_s/\delta_s=s/2\to0$. Direct substitution gives the exact
pointwise identity
\begin{equation}\label{mf:eq:exactcore}
\ell_s^2H_s(q_\pm+\ell_s\xi)
=\frac{8|\pm1+(s/2)\xi|^2}
{\bigl(1+|\pm\xi+(s/4)\xi^2|^2\bigr)^2}.
\end{equation}
On each fixed $B_R$, this converges in $C^2$ to
$8/(1+|\xi|^2)^2$ with error $O_R(s)$.

For a sufficiently small fixed $c>0$, factorization of $z^2-s$
implies
\begin{equation}\label{mf:eq:pointwise}
H_s(z)\asymp
\ell_s^{-2}\left(1+\frac{|z-q_\pm|^2}{\ell_s^2}\right)^{-2}
\qquad (|z-q_\pm|\le c\sqrt{s}),
\end{equation}
where the comparison constants are independent of small $s$.
We shall also use
\begin{equation}\label{mf:eq:supH}
\|H_s\|_{L^\infty(B_1)}\le64s^{-3},\qquad 0<s<1.
\end{equation}
Indeed, for $|z|^2\le2s$ this follows from \eqref{mf:eq:seed} and the
lower bound $s^4$ for $s^4+|z^2-s|^2$. For $|z|^2>2s$, the bound
$|z^2-s|\ge |z|^2/2$ gives $H_s\le512s^4/|z|^6\le64s$.

\subsection{Selection of nondegenerate seeds}

Fix a H\"older exponent $\sigma_H\in(0,1)$ and let
\[
X=\{\phi\in C^{2,\sigma_H}(\overline{B_1}):
\phi|_{\partial B_1}=0\},\qquad
Y=C^{0,\sigma_H}(\overline{B_1}).
\]
Define
\begin{equation}\label{mf:eq:linear}
L_s=-\Delta-H_s:X\longrightarrow Y.
\end{equation}

\begin{lemma}\label{mf:lem:generic}
There is a discrete subset $\mathcal E_D\subset(0,\infty)$ such that $L_s$ is
invertible for every $s\in(0,\infty)\setminus \mathcal E_D$.
\end{lemma}
\begin{proof}
Writing $r^2=x^2+y^2$, we have
\(
H_s(x,y)=\frac{32s^4r^2}
{\bigl(s^4+s^2-2s(x^2-y^2)+r^4\bigr)^2}.
\)
For positive real $s$, the denominator is nonzero on the closed disk.
Small complex disks about the positive parameters keep the denominator
uniformly separated from zero on $\overline{B_1}$. Their union is a
connected neighborhood of $(0,\infty)$ on which $H_s$ is holomorphic.

Let $A=-\Delta_D$ on $L^2(B_1)$ and let $M_{H_s}$ denote
multiplication by $H_s$. Then $A^{-1}M_{H_s}$ is a holomorphic family
of compact operators on $L^2(B_1)$. For a concrete invertible parameter, $H_2\le2$ on $B_1$ and
\(\lambda_1(B_1)\ge\lambda_1((-1,1)^2)=\pi^2/2>2.\)
Hence $I-A^{-1}M_{H_2}$ is invertible. The analytic Fredholm theorem
\cite[Theorem~5.3.1]{FredholmNotes} makes the exceptional set discrete;
elliptic regularity and the Fredholm alternative give the same
invertibility between $X$ and $Y$.
\end{proof}

Choose
\begin{equation}\label{mf:eq:selects}
s_k\in\left(\frac1{k+2},\frac1{k+1}\right)\setminus \mathcal E_D.
\end{equation}
This is possible because $\mathcal E_D$ contains no interval.
Then $s_k\downarrow0$ and each $L_{s_k}$ is invertible;
no uniform bound for $L_{s_k}^{-1}$ is required.

\subsection{The nonlinear perturbation}

Fix $s\notin \mathcal E_D$, put $\tau=\eps^2$, and seek a solution in the form
\begin{equation}\label{mf:eq:ansatz}
u=\log\tau+2\log|z|+v_s+\phi,\qquad \phi\in X.
\end{equation}
Since $\Delta(2\log|z|)=4\pi\delta_0$, equations
\eqref{mf:eq:seedPDE} and \eqref{mf:eq:ansatz} reduce \eqref{mf:eq:CS} exactly to
\begin{equation}\label{mf:eq:nonlinear}
\mathcal F_s(\phi,\tau):=
-\Delta\phi-H_s(e^\phi-1)+\tau H_s^2e^{2\phi}=0.
\end{equation}
The map $\mathcal F_s:X\times\R\to Y$ is smooth, and
\(\mathcal F_s(0,0)=0,\qquad D_\phi\mathcal F_s(0,0)=L_s.\)
The Banach-space implicit function theorem supplies a solution
$\phi_s(\tau)$ for $|\tau|$ sufficiently small, with $\phi_s(0)=0$.
Moreover,
\begin{equation}\label{mf:eq:phiExpansion}
\phi_s(\tau)=-\tau L_s^{-1}(H_s^2)+O_s(\tau^2)
\qquad\text{in }C^{2,\sigma_H}(\overline{B_1}).
\end{equation}
The allowed interval and the constants may depend on $s$.

For $s_k$ chosen in \eqref{mf:eq:selects}, take $\tau_k>0$ in its
implicit-function interval so small that
\begin{equation}\label{mf:eq:diagonal}
\|\phi_{s_k}(\tau_k)\|_{C^{2,\sigma_H}}\le s_k,
\qquad \tau_k\le s_k^4,
\qquad \tau_k<\tfrac14\tau_{k-1}\quad(k>1).
\end{equation}
Set $\eps_k=\sqrt{\tau_k}$, $\phi_k=\phi_{s_k}(\tau_k)$, and
\begin{equation}\label{mf:eq:solution}
u_k=2\log\eps_k+2\log|z|+v_{s_k}+\phi_k.
\end{equation}
These solve \eqref{mf:eq:CS} exactly, with smooth regular parts by
elliptic bootstrapping. Since $\phi_k=0$ on $\partial B_1$, they have
the traces \eqref{mf:eq:boundary}, whose oscillations are $O(s_k)$.

\subsection{Blow-up type, core estimates, and concentration}

By \eqref{mf:eq:supH} and \eqref{mf:eq:diagonal},
\begin{equation}\label{mf:eq:MF}
e^{u_k}=\tau_kH_{s_k}e^{\phi_k}
\le64e^{s_k}\tau_ks_k^{-3}
\le64e^{s_k}s_k\longrightarrow0.
\end{equation}
Thus $\sup_{B_1}u_k\to-\infty$. After discarding finitely many terms,
$u_k<0$ and $D_k\ge0$.

The regularized functions are $w_k=v_{s_k}+\phi_k$. At
$q_{\pm,k}=\pm\sqrt{s_k}$, we have
\begin{align}\label{mf:eq:failure}
w_k(q_{\pm,k})+4\log|q_{\pm,k}|
&=\log32-4\log s_k+2\log s_k+\phi_k(q_{\pm,k})\notag\\
&=2\log(1/s_k)+\log32+o(1)\longrightarrow+\infty.
\end{align}
Since $w_k\to-\infty$ locally away from zero and
$w_k(q_{\pm,k})\to+\infty$, the origin is its unique blow-up point;
\eqref{mf:eq:failure} proves non-simplicity.

Equations \eqref{mf:eq:exactcore} and \eqref{mf:eq:diagonal} yield
\begin{equation}\label{mf:eq:localprofile}
u_k(q_{\pm,k}+\ell_k\xi)-u_k(q_{\pm,k})
=-2\log(1+|\xi|^2)+O_R(s_k)
\quad\text{in }C^2(B_R)
\end{equation}
for each fixed $R$, where $\ell_k=s_k^{3/2}/2$.
Indeed, normalize \eqref{mf:eq:exactcore} by its value at zero,
take logarithms, and use $\|\phi_k\|_{C^2}\le s_k$; its numerator
stays positive on fixed $B_R$. The nondegenerate limiting maximum
gives maxima at $q_{\pm,k}+o(\ell_k)$, with separation
$2\sqrt{s_k}(1+o(1))$ and relative width $\ell_k/\sqrt{s_k}=s_k/2\to0$.

The exact nonlinear source is
\begin{equation}\label{mf:eq:D}
D_k=H_{s_k}e^{\phi_k}
\bigl(1-\tau_kH_{s_k}e^{\phi_k}\bigr).
\end{equation}
Using \eqref{mf:eq:totalmass}, \eqref{mf:eq:supH}, and \eqref{mf:eq:diagonal},
we obtain
\begin{align}\label{mf:eq:L1}
\|D_k-H_{s_k}\|_{L^1(B_1)}
&\le C\|\phi_k\|_\infty\int_{B_1}H_{s_k}\dd
 +C\tau_k\|H_{s_k}\|_\infty\int_{B_1}H_{s_k}\dd\notag\\
&\le Cs_k\longrightarrow0.
\end{align}

For mass exhaustion, fix $c\in(0,1/4)$. Outside
$B_{c\sqrt{s}}(\sqrt{s})\cup B_{c\sqrt{s}}(-\sqrt{s})$, write
$y=z/\sqrt{s}$. Both $|y-1|$ and $|y+1|$ are at least $c$, so
$|y^2-1|\ge c^2$ and $|f_s(z)|\ge c^2/s$. The degree-two
change of variables in \eqref{mf:eq:seedPDE} bounds the mass of this
region by
\[
2\int_{|w|\ge c^2/s}\frac8{(1+|w|^2)^2}\,dw=O(s^2).
\]
Symmetry and \eqref{mf:eq:totalmass} give mass $8\pi+O(s^2)$ in
each of the two core disks. These disks lie in $B_1$ for small $s$.
Combining this fact with \eqref{mf:eq:L1} proves
\(D_k\dd\rightharpoonup16\pi\delta_0.\)
Furthermore, \eqref{mf:eq:exactcore} and \eqref{mf:eq:L1} give
\[
\lim_{R\to\infty}\lim_{k\to\infty}
\int_{B_{R\ell_k}(q_{\pm,k})}D_k\dd
=\int_{\R^2}\frac8{(1+|\xi|^2)^2}\,d\xi=8\pi.
\]
The pointwise comparison \eqref{mf:eq:pointwise} remains valid for
$D_k$ on the core disks, up to uniform multiplicative constants,
by \eqref{mf:eq:MF} and \eqref{mf:eq:D}. This completes the proof of
Theorem~\ref{mf:thm:main}.

The corresponding boundary flux follows by integrating \eqref{mf:eq:CS}:
\begin{equation}\label{mf:eq:flux}
\int_{\partial B_1}\partial_{\boldsymbol n}u_k\,ds
=4\pi-\int_{B_1}D_k\dd\longrightarrow-12\pi.
\end{equation}
Thus one strength-one vortex can support mass $16\pi$ on the disk.
On a closed surface, total mass $4\pi N_{\mathrm{tot}}$ would instead
require $N_{\mathrm{tot}}=4$ for an exhausting mean-field pair.
 

\appendix
\section{Proofs and moment formulas for finite-height cores}
\label{app:pointwise-proofs}

This appendix proves Theorem~\ref{prop:refined-profile} and develops
its moment-sampling consequences. We retain the finite-height
assumption and notation of Section~\ref{sec:pointwise}.

\subsection{Local normalization and initial estimates}
\label{app:core-normalization}

Under Assumption~\ref{ass:profile}, retain the first-scale equation,
centers, and actual masses from Section~\ref{sec:cluster}. We use the
logarithmic comparison of Lemma~\ref{lem:isolated-core-log} and the
normalized-kernel rigidity of Lemma~\ref{thm:linearized-classification}.

Choose the local disk inside a slightly larger isolated-core
neighborhood. On a fixed annulus containing its boundary, the Green
representation and exterior oscillation estimate bound
$v_\eps-\fint_{\partial B_r(q_{t,\eps})}v_\eps$ uniformly.
The source is bounded there by \eqref{eq:isolated-core-log};
interior elliptic estimates and differentiation of the equation
therefore give uniform $C^{4,\sigma}$ bounds on a smaller annulus,
for any fixed $0<\sigma<1$.
For each $t$, let $\widetilde{\mathfrak h}_{t,\eps}$ be the harmonic function
satisfying
\begin{equation}
  \begin{cases}
  \Delta\widetilde{\mathfrak h}_{t,\eps}=0
  &\text{in} \ B_r(q_{t,\eps}),\\
  \displaystyle
  \widetilde{\mathfrak h}_{t,\eps}
  =v_\eps-\fint_{\partial B_r(q_{t,\eps})}v_\eps\,dS
  &\text{on}\ \partial B_r(q_{t,\eps}).
  \end{cases}
  \label{eq:harmonic-correction}
\end{equation}
Let $q_{t,\eps}^*$ be the maximum point of
$v_\eps-\widetilde{\mathfrak h}_{t,\eps}$ in $B_r(q_{t,\eps})$.
The bounded harmonic correction and \eqref{eq:isolated-core-log}
give $q_{t,\eps}^*-q_{t,\eps}=O(\etae)$, so this point remains a
uniform positive distance from the boundary. Subtract its value at
$q_{t,\eps}^*$ to normalize
\begin{equation}
  \mathfrak h_{t,\eps}
  :=\widetilde{\mathfrak h}_{t,\eps}
  -\widetilde{\mathfrak h}_{t,\eps}(q_{t,\eps}^*),
  \qquad
  \mathfrak h_{t,\eps}(q_{t,\eps}^*)=0.
  \label{eq:phi-normalization}
\end{equation}
Thus $\mathfrak u_{t,\eps}:=v_\eps-\mathfrak h_{t,\eps}$ is constant on the boundary.
Boundary Schauder estimates also give
$\|\mathfrak h_{t,\eps}\|_{C^{4,\sigma}(\overline{B_r(q_{t,\eps})})}\le C$.
All subsequent coefficient Taylor estimates consequently hold on the
full closed local disk.

Set
\begin{equation}\label{eq:H-alpha-local}
 \begin{aligned}
 H_{t,\eps}&:=a_\eps e^{\mathfrak h_{t,\eps}},\quad
 s_{t,\eps}:=\etae e^{-\varphi_i(\de q_{t,\eps}^*)/2},\quad
\alpha_{t,\eps}:=v_\eps(q_{t,\eps}^*)+2\log\etae
                   +\log a_\eps(q_{t,\eps}^*).
 \end{aligned}
\end{equation}
The two maximum properties and the Lipschitz bound on $\mathfrak h_{t,\eps}$ give
\[
 0\le v_\eps(q_{t,\eps})-v_\eps(q_{t,\eps}^*)
 \le C|q_{t,\eps}-q_{t,\eps}^*|=O(\etae).
\]
Together with \eqref{eq:a-uniform}, this yields
$\alpha_{t,\eps}=\vartheta_{t,\eps}
+\log a_\eps(q_{t,\eps})+O(\etae)\to\beta$.
Choosing $I$ in \eqref{eq:profile-hyp} with $\beta$ in its interior,
we have $\alpha_{t,\eps}\in I$ for small $\eps$ and
$m_{t,\eps}:=m(\alpha_{t,\eps})>m_*>4$. Define {
\begin{align}
\widehat v_{t,\eps}(z)
  &:={\mathfrak u}_{t,\eps}(q_{t,\eps}^*+s_{t,\eps}z)
    +2\log\etae+\log H_{t,\eps}(q_{t,\eps}^*),
    \quad
  A_{t,\eps}(z)
   :=\frac{H_{t,\eps}(q_{t,\eps}^*+s_{t,\eps}z)}
  {H_{t,\eps}(q_{t,\eps}^*)}.
  \label{eq:A-local}
\end{align}
Writing $\ell_{t,\eps}=\log A_{t,\eps}$ and
$\Gamma_{t,\eps}(z)=\varphi_i(\de(q_{t,\eps}^*+s_{t,\eps}z))
-\varphi_i(\de q_{t,\eps}^*)$, we obtain
\begin{equation}\label{eq:local-rescaled-equation}
 \Delta\widehat v_{t,\eps}+e^{\Gamma_{t,\eps}}
 F(\widehat v_{t,\eps}+\ell_{t,\eps})=0,\qquad
 \widehat v_{t,\eps}(0)=\alpha_{t,\eps},\quad
 \nabla\widehat v_{t,\eps}(0)=0.
\end{equation}

By local compactness and the regular-profile classification recalled
in \cite[pp.~840--841]{CKL}, with
$U_{t,\eps}:=U_{\alpha_{t,\eps}}$,
$\widehat v_{t,\eps}-U_{t,\eps}\to0$ in $C^2_{\loc}(\R^2)$.
Positivity of the source gives
$\liminf\widetilde m_{t,\eps}\ge m(\beta)>4$ along a subsequence
with $\alpha_{t,\eps}\to\beta$.
By the strict choice in \eqref{eq:profile-hyp}, both the actual local
slopes $\widetilde m_{t,\eps}$ and the profile slopes exceed the fixed
$m_*$ for all sufficiently small $\eps$.
The isolated-core estimate \eqref{eq:isolated-core-log}, used
with its actual slope, now gives}
\begin{equation}
  e^{\widehat v_{t,\eps}(z)}
  +|F'(\widehat v_{t,\eps}(z))|
  \le C(1+|z|)^{-m_*}
  \label{eq:scaled-source-decay}
\end{equation}
throughout the rescaled local disk. Local convergence together with
this integrable tail bound gives
$\widetilde m_{t,\eps}=m_{t,\eps}+o(1)$.

Put $W_{t,\eps}:=\widehat v_{t,\eps}-U_{t,\eps}$.
The rescaled disk is
\begin{equation}\label{eq:rescaled-core-disk}
 D_{t,\eps}=B_{R_{t,\eps}}(\zeta_{t,\eps}),\qquad
 R_{t,\eps}=\frac r{s_{t,\eps}},\qquad
 \zeta_{t,\eps}=\frac{q_{t,\eps}-q_{t,\eps}^*}{s_{t,\eps}}.
\end{equation}
The symbol $\mathfrak u_{t,\eps}$ always denotes this local harmonic
regularization; $u_\eps$ remains the original physical solution.
We suppress fixed-core indices: $\alpha=\alpha_{t,\eps}$,
$U=U_{t,\eps}$, $m=m_{t,\eps}$, while
$s_\eps,D_\eps,R_\eps,\zeta_\eps$ retain only the coupling index.

The local logarithmic coefficient has a small trace on the first scale.
Indeed, $\log H_{t,\eps}=\log a_\eps+\mathfrak h_{t,\eps}$
is harmonic on $B_r(q_{t,\eps})$, which avoids the vortex.
Thus $\operatorname{tr}\mathsf B_{t,\eps}=0$ exactly, and only the
conformal factor contributes to \(
 \left|
 \left.\Delta_y\log\bigl(e^{\varphi_i(\de y)}H_{t,\eps}(y)\bigr)
 \right|_{y=q_{t,\eps}^*}
 \right|\le C\de^2.\)

\subsection{Preliminary center and first-moment estimates}
\label{app:preliminary-estimates}

\begin{lemma}[Preliminary comparison of the two centers]
\label{lem:preliminary-center}
With the notation above,
\begin{equation}
  |q_{t,\eps}^*-q_{t,\eps}|\le C\etae^2.
  \label{eq:preliminary-center}
\end{equation}
\end{lemma}

\begin{proof}
The isolated-core comparison and strict radial monotonicity give
$\zeta_{t,\eps}\to0$. Since $q_{t,\eps}$ is a critical point of
$v_\eps=\mathfrak u_{t,\eps}+\mathfrak h_{t,\eps}$,
\begin{equation}
  0={s_{t,\eps}\nabla v_\eps(q_{t,\eps})
  =\nabla\widehat v_{t,\eps}(\zeta_{t,\eps})+s_{t,\eps}\nabla\mathfrak h_{t,\eps}(q_{t,\eps}).
  \label{eq:preliminary-center-gradient}}
\end{equation}
Taylor's formula between $0$ and { $\zeta_{t,\eps}$,} the convergence
{ $\widehat v_{t,\eps}\to U_\alpha$} in $C^2_{\loc}$, and
$D^2U_\alpha(0)=-\kappa_\alpha I_2/2$ show that the Hessian in this
formula is uniformly invertible.  Since { $\nabla\mathfrak h_{t,\eps}$} is uniformly bounded,
\eqref{eq:preliminary-center-gradient} gives { $|\zeta_{t,\eps}|\le C\etae$,} which is
\eqref{eq:preliminary-center}.
\end{proof}

\begin{lemma}[Preliminary weighted estimate]
\label{lem:preliminary-pointwise}
Under the assumptions of Theorem~\ref{prop:refined-profile},
\begin{equation}\label{eq:preliminary-pointwise}
 |W_{t,\eps}(z)|\le C\etae(1+|z|)^\sigma,
 \qquad |\nabla W_{t,\eps}(z)|\le C\etae
 \quad\left(|z|\le\frac{3r}{4s_{t,\eps}}\right).
\end{equation}
In fact the estimate for $W_{t,\eps}$ holds on the full rescaled disk.
\end{lemma}

\begin{proof}
Suppress the core indices in $U$, $W$, $\ell$, and $\Gamma$.
By Lemma~\ref{lem:preliminary-center}, $\zeta_\eps=O(\etae)$.
The coefficient bounds on the full local disk give
\begin{equation}\label{eq:A-crude}
 |\ell(z)|+|\Gamma(z)|\le C\etae|z|\le C
 \qquad(z\in D_\eps).
\end{equation}
Subtract the equations by exact linearization:
\begin{equation}\label{eq:W-equation-exact}
 \Delta W+c_\eps W=\mathfrak f_\eps,
 \qquad
 \begin{aligned}
 c_\eps(z)&=e^{\Gamma(z)}\int_0^1
 F'(U(z)+\ell(z)+\theta W(z))\,d\theta,\\
 \mathfrak f_\eps(z)&=F(U(z))-e^{\Gamma(z)}F(U(z)+\ell(z)).
 \end{aligned}
\end{equation}
The exponential decay \eqref{eq:scaled-source-decay} applies at both
endpoints $U+\ell$ and $\widehat v+\ell$. Their geometric interpolation
therefore gives
\begin{equation}\label{eq:exact-linearization-bounds}
 |c_\eps(z)|\le C(1+|z|)^{-m_*},\qquad
 |\mathfrak f_\eps(z)|\le C\etae(1+|z|)^{1-m_*}.
\end{equation}
Also $c_\eps\to F'(U_\beta)$ locally along any subsequence for which
$\alpha_{t,\eps}\to\beta$. No bound on $\|W\|_\infty$ has been used.

Let $h_\eps$ be the harmonic extension of $W|_{\partial D_\eps}$.
Since $\widehat v$ is constant there, radiality and
$|U'(r)|\le C/r$ give
\begin{equation}\label{eq:harmonic-boundary-error}
 \operatorname{osc}_{D_\eps}h_\eps
 \le\operatorname{osc}_{\partial D_\eps}U
 \le C|\zeta_\eps|/R_\eps\le C\etae^2.
\end{equation}
For the positive Dirichlet Green function of $-\Delta$,
$-\Delta W=c_\eps W-\mathfrak f_\eps$. Since $W(0)=0$, subtraction at the
origin yields
\begin{equation}\label{eq:green-preliminary-contradiction}
 W(z)=h_\eps(z)-h_\eps(0)+\int_{D_\eps}
 [G_\eps(z,y)-G_\eps(0,y)]
 [c_\eps(y)W(y)-\mathfrak f_\eps(y)]\,dy.
\end{equation}
Lemma~\ref{lem:weighted-green} also applies to $D_\eps$: translate
the disk and subtract its values at the translated origin; the weights
are uniformly equivalent since $\zeta_\eps=O(\etae)$.

Suppose that $\mathcal N_\eps:=\max_{\overline D_\eps}
|W(z)|/(1+|z|)^\sigma$ satisfies
$\Lambda_\eps:=\mathcal N_\eps/\etae\to\infty$, and set
$\widehat W_\eps=W/\mathcal N_\eps$. Then
$|\widehat W_\eps(z)|\le(1+|z|)^\sigma$ and
$\widehat W_\eps(0)=\nabla\widehat W_\eps(0)=0$.
Local elliptic compactness and \eqref{eq:exact-linearization-bounds}
show that every local limit solves $L_\beta\widehat W=0$ with these
normalizations. Lemma~\ref{thm:linearized-classification} makes
that limit zero. A maximizer of the weighted norm must consequently
escape every fixed disk.

Divide \eqref{eq:green-preliminary-contradiction} at such a maximizer
$z_\eps$ by $\mathcal N_\eps(1+|z_\eps|)^\sigma$. The harmonic term is
$O(\etae/\Lambda_\eps)$, and the forcing potential is
$O(\Lambda_\eps^{-1})$ by Lemma~\ref{lem:weighted-green}. For the
remaining potential first restrict to a fixed $B_R$, where
$\widehat W_\eps\to0$. On its complement choose
$0<\rho<m_*-2-\sigma$ and use
\[
 |c_\eps\widehat W_\eps|
 \le C(1+|y|)^{-m_*+\sigma}
 \le C R^{-(m_*-2-\sigma-\rho)}(1+|y|)^{-2-\rho}.
\]
The weighted Green estimate makes the tail arbitrarily small,
uniformly in $\eps$. The resulting zero limit contradicts the unit
weighted value at $z_\eps$. Hence $\mathcal N_\eps\le C\etae$ on the full
disk. Differentiating the Green formula on $|z|\le3r/(4s_\eps)$ gives
$|\nabla W|\le C\etae$: the source is bounded by
$C\etae(1+|z|)^{-m_*+1}$, and the harmonic gradient is $O(\etae^3)$.

In particular $\|W\|_{L^\infty(D_\eps)}
\le C\etae^{1-\sigma}\to0$, which justifies the quadratic Taylor
expansions used below.
\end{proof}

\begin{lemma}[First moment on a tangent circle]\label{lem:first-moment}
Suppress the core indices, including $\mathfrak u=\mathfrak u_{t,\eps}$.
Put $\mathbf{d}=q^*-q$, $r_\eps=r-|\mathbf{d}|$, and
$\boldsymbol\omega(\theta)=(\cos\theta,\sin\theta)$.
Then
\begin{equation}\label{eq:first-moment-coarse}
 \frac1{2\pi}\int_0^{2\pi}\mathfrak u(q^*+r_\eps \boldsymbol\omega(\theta))\boldsymbol\omega(\theta)\,d\theta
 =-\frac{m}{2r}\mathbf{d}+O(\etae^2).
\end{equation}
\end{lemma}

\begin{proof}
Suppress $t,\eps$, translate $q$ to zero, and write $q^*=\mathbf{d}$. Since $\mathfrak u$
is constant on $\partial B_r$, its Green representation is
\(
 \mathfrak u(x)=C+\int_{B_r}G_r(x,y)f_\eps(y)\,dy.
\)
Here
$G_r(x,y)=-(2\pi)^{-1}\log|x-y|
 +(2\pi)^{-1}\log(|r^2-\bar yx|/r)=:G_r^{\rm sing}+G_r^{\rm reg}$.
Under $y=\mathbf{d}+s_\eps z$, the source becomes
\begin{equation}\label{eq:correct-source-change-variable}
 f_\eps(y)\,dy=\mathcal F_\eps(z)\,dz,\qquad
 \mathcal F_\eps:=e^\Gamma F(U+W+\ell).
\end{equation}
Restoring the core index, we write $\mathcal F_{t,\eps}=\mathcal F_\eps$.
Lemma~\ref{lem:preliminary-pointwise} and the coefficient expansion give
\begin{equation}\label{eq:source-moment-error}
 |\mathcal F_\eps-F(U)|\le C\etae(1+|z|)^{1+\sigma-m_*},
 \qquad
 \int(1+|z|)|\mathcal F_\eps-F(U)|\,dz\le C\etae.
\end{equation}
The integrals are over the rescaled disk. The unweighted source mass
in the far tail $|z|\ge c/\etae$ is
$O(\etae^{m_*-2})=o(\etae^2)$.

For $\Pi_\eps f=(2\pi)^{-1}
\int_0^{2\pi}f(\mathbf{d}+r_\eps \boldsymbol\omega(\theta))\boldsymbol\omega(\theta)\,d\theta$, the identity
\[
 \frac1{2\pi}\int_0^{2\pi}
 \log|r_\eps \boldsymbol\omega(\theta)-\zeta|\,\boldsymbol\omega(\theta)\,d\theta
 =-\frac12
 \begin{cases}
 \zeta/r_\eps,&|\zeta|<r_\eps,\\
 r_\eps\zeta/|\zeta|^2,&|\zeta|>r_\eps
 \end{cases}
\]
shows that the radial leading source has zero singular first mode on
$|z|<r_\eps/s_\eps$; its remaining tail is $o(\etae^2)$.
The source error is controlled by \eqref{eq:source-moment-error}, giving
\(
 \int_{B_r}\Pi_\eps[G_r^{\rm sing}(\,\cdot\,,y)]f_\eps(y)\,dy
 =O(\etae^2).
\)
The regular projected kernel has the exact complex formula
\begin{equation}\label{eq:regular-projected-kernel}
 K_\eps(y):=\Pi_\eps[G_r^{\rm reg}(\,\cdot\,,y)]
 =-\frac{r_\eps y}{4\pi(r^2-\overline{\mathbf{d}} y)}.
\end{equation}
It is harmonic with uniformly bounded derivatives on $B_r$.
The radial mean-value identity, with the same tail bound, gives
\begin{equation}\label{eq:regular-radial-reduction}
 \int K_\eps(\mathbf{d}+s_\eps z)F(U(z))\,dz
 =2\pi mK_\eps(\mathbf{d})+o(\etae^2).
\end{equation}
\begin{equation}\label{eq:regular-leading-first-mode}
 2\pi mK_\eps(\mathbf{d})=-\frac{m \mathbf{d}}{2(r+|\mathbf{d}|)}
 =-\frac{m}{2r}\mathbf{d}+O(\etae^4).
\end{equation}
Because $K_\eps(\mathbf{d})=O(|\mathbf{d}|)$ and $K_\eps$ is uniformly Lipschitz near
$\mathbf{d}$, replacing $F(U)$ by $\mathcal F_\eps$ costs $O(\etae^2)$.
Combining the two parts proves \eqref{eq:first-moment-coarse}.
\end{proof}

\subsection{Proof of the quantitative approximation theorem}
\label{app:pointwise-proof}

\begin{proof}[Proof of Theorem~\ref{prop:refined-profile}]
We suppress the core indices and put $\kappa=F(\alpha)$,
$\mathbf{d}=q^*-q$, and
{ 
\begin{equation}
  \mathbf{b}:=\nabla\log H(q^*),
  \qquad
  \mathbf{g}:=\de\nabla_y\varphi_i(y)\Big|_{y=\de q^*},
  \qquad
  \mathbf{b}^{\rm eff}:=\mathbf{b}+\frac m4\mathbf{g},
  \label{eq:b-g-effective}
\end{equation}
}
The conformal normalization gives
\begin{equation}
  \mathbf{g}=O(\de^2),
  \qquad
  s_\eps=\etae(1+O(\de^2)).
  \label{eq:s-g-size}
\end{equation}

Lemma~\ref{lem:preliminary-center} gives $\mathbf{d}=O(\etae^2)$ and
$\zeta_\eps=-\mathbf{d}/s_\eps=O(\etae)$. The maximum condition is
\begin{equation}
  0
  =
  s_\eps\nabla v_\eps(q)
  =
  \nabla\widehat v(\zeta_\eps)
  +s_\eps\nabla\mathfrak h(q).
  \label{eq:max-condition-scaled}
\end{equation}

\needspace{5\baselineskip}
\medskip
\noindent
{ 
\emph{Step 1: the conformal first Fourier correction.}
}

Taylor expansion gives
{ 
\begin{align}
  \log A(z)
  &=
  s_\eps \mathbf{b}\cdot z
  +O(s_\eps^2|z|^2),
  \ \
  \varphi_i(\de(q^*+s_\eps z))
  -\varphi_i(\de q^*)
  =
  s_\eps \mathbf{g}\cdot z
  +O(\de^2s_\eps^2|z|^2).
  \label{eq:conformal-first}
\end{align}
}
By Lemma~\ref{lem:preliminary-pointwise},
{ 
\begin{align}
  L_\alpha W
  ={}&
  -s_\eps(\mathbf{b}\cdot z)F'(U_\alpha)
  -s_\eps(\mathbf{g}\cdot z)F(U_\alpha)
  +E_\eps^{(2)},
  \label{eq:W-first-decomposition}
\end{align}
}
where \begin{equation}
\begin{aligned}
  |E_\eps^{(2)}(z)|
  &\le 
  c\left(\eta_\eps^2(1+\de^2)
  (1+|z|)^{2-m_*}
 +
  (1+|z|)^{-m_*}|W(z)|^2\right)  \le C \eta_\eps^2
  (1+|z|)^{2-m_*}.
  \label{eq:E2-crude}
\end{aligned}\end{equation}

For the coefficient first mode set
\begin{equation}
  Z_{\mathbf{b}}(z)
  :=
  -s_\eps \mathbf{b}\cdot z
  -\frac{2s_\eps}{\kappa}
   \mathbf{b}\cdot\nabla U_\alpha(z).
  \label{eq:explicit-first-mode}
\end{equation}
Then
\(L_\alpha Z_{\mathbf{b}} = -s_\eps(\mathbf{b}\cdot z)F'(U_\alpha).\)

{ 
For the conformal first mode, let $Y_\alpha$ be the normalized regular
solution of
\begin{equation}
  Y_\alpha''
  +\frac1{\varrho}Y_\alpha'
  -\frac1{\varrho^2}Y_\alpha
  +F'(U_\alpha)Y_\alpha
  =
  -\varrho F(U_\alpha),
  \qquad
  Y_\alpha(0)=Y_\alpha'(0)=0.
  \label{eq:Y-first-mode}
\end{equation}}
The normalized solution has the explicit form
\begin{equation}\label{eq:Y-explicit}
 Y_\alpha(\varrho)=\frac{\varrho^2}{4}U_\alpha'(\varrho).
\end{equation}
Indeed, let $L_{\alpha,1}=\partial_\varrho^2+
\varrho^{-1}\partial_\varrho-\varrho^{-2}+F'(U_\alpha)$.
Since $L_{\alpha,1}U_\alpha'=0$, the radial profile equation gives
\[
 L_{\alpha,1}(\varrho^2U_\alpha')
 =4(\varrho U_\alpha''+U_\alpha')
 =-4\varrho F(U_\alpha).
\]
Formula~\eqref{eq:Y-explicit} has the required normalization at zero;
uniqueness follows because the regular homogeneous first mode is a
multiple of $U_\alpha'$, whose derivative at zero is $-\kappa/2\ne0$.
It also gives, with differentiated expansions,
\[
 \begin{aligned}
 Y_\alpha(\varrho)&=-\kappa\varrho^3/8+O(\varrho^5)
 &&(\varrho\to0),\quad\
 Y_\alpha(\varrho)&=-\frac m4\varrho+O(\varrho^{3-m_*})
 &&(\varrho\to\infty).
 \end{aligned}
\]
Consequently
\begin{equation}\label{eq:Q-first-mode}
 \begin{aligned}
 Q_\alpha(\varrho)
 &:=Y_\alpha(\varrho)+\frac m4\varrho
       +\frac{m}{2\kappa}U_\alpha'(\varrho) =\left(\frac{\varrho^2}{4}+\frac{m}{2\kappa}\right)
       U_\alpha'(\varrho)+\frac m4\varrho
 \end{aligned}
\end{equation}
satisfies $Q_\alpha=O(\varrho^3)$, $Q_\alpha'=O(\varrho^2)$ at zero
and $Q_\alpha=O(\varrho^{-1})$, $Q_\alpha'=O(\varrho^{-2})$ at infinity.

Set
{ 
\begin{equation}
  \Xi_1(z)
  :=
  s_\eps Q_\alpha(|z|)
  \mathbf{g}\cdot\frac z{|z|}.
  \label{eq:Xi1-definition}
\end{equation}
}
Then
\begin{equation}
  |\Xi_1(z)|+|\nabla\Xi_1(z)|
  \le C\etae\de^2.
  \label{eq:Xi1-bound}
\end{equation}

With $Z_{\mathbf g}(z):=s_\eps Y_\alpha(|z|)\mathbf g\cdot z/|z|$,
the definitions of $\mathbf b^{\rm eff}$ and $Q_\alpha$ give
\begin{align}
  Z_{\mathbf{b}}(z)+Z_{\mathbf{g}}(z)
  ={}&
  -s_\eps \mathbf{b}^{\rm eff}\cdot z
  -\frac{2s_\eps}{\kappa}
   \mathbf{b}^{\rm eff}\cdot\nabla U_\alpha(z)
  +\Xi_1(z).
  \label{eq:first-mode-effective-decomposition}
\end{align}

Let \(\widetilde W:=W-\Xi_1\).
Then \eqref{eq:W-first-decomposition} implies that
{ 
\begin{equation}
  L_\alpha\widetilde W
  =
  -s_\eps(\mathbf{b}^{\rm eff}\cdot z)F'(U_\alpha)
  + E_\eps^{(2)},
  \label{eq:Wtilde-equation}
\end{equation}
}
with the same remainder \eqref{eq:E2-crude}.

\medskip
\noindent{ 
\emph{Step 2: a refined weighted estimate.}
}
We claim that
\begin{equation}
  |\widetilde W(z)|
  \le
  Cs_\eps\bigl(|\mathbf{b}^{\rm eff}|+s_\eps\bigr)(1+|z|)^\sigma, 
  \qquad
  |z|\le\frac{3r}{4s_\eps}.
  \label{eq:Wtilde-refined}
\end{equation}
To prove this, apply the full-disk normalization and Green argument
of Lemma~\ref{lem:preliminary-pointwise} to
$\widetilde W=W-\Xi_1$, with comparison scale
$\mathfrak n_\eps=s_\eps(|\mathbf{b}^{\rm eff}|+s_\eps)$. The source in
\eqref{eq:Wtilde-equation} divided by $\mathfrak n_\eps$ is bounded by
$C(1+|z|)^{2-m_*}$, and the boundary oscillation is
$O(\etae^2)+O(s_\eps^2\de^2)=O(\mathfrak n_\eps)$.
Consequently the normalized forcing and boundary terms vanish in a
contradiction sequence. Local limits vanish by
Lemma~\ref{thm:linearized-classification}, and the Green tails are
uniformly small since $m_*>4+\sigma$. This proves
\eqref{eq:Wtilde-refined}; its weighted bound holds on the full disk
$D_\eps$.

We next claim that
\begin{equation}
  |\mathbf{b}^{\rm eff}|\le Cs_\eps.
  \label{eq:beff-O-s}
\end{equation}
Otherwise, along a subsequence,
$|\mathbf{b}^{\rm eff}|/s_\eps\to\infty$,
$\mathbf{b}^{\rm eff}/|\mathbf{b}^{\rm eff}|\to\mathbf e_*$, $|\mathbf e_*|=1$, and
$\alpha_\eps\to\alpha_*$. The functions
$\widetilde W/(s_\eps|\mathbf{b}^{\rm eff}|)$ converge locally to $\Psi$ with
\begin{equation}\label{eq:Psi-effective-limit}
 L_{\alpha_*}\Psi=-(\mathbf e_*\cdot z)F'(U_{\alpha_*}),\qquad
 \Psi(0)=0,\quad\nabla\Psi(0)=0,\quad
 |\Psi(z)|\le C(1+|z|)^\sigma.
\end{equation}
Rotate so that $\mathbf e_*=(1,0)$. The cosine first mode $\Psi_1$ satisfies
the forced $k=1$ equation, so $\Psi_1(\varrho)+\varrho$ solves its homogeneous
equation. Regularity at zero implies
$\Psi_1(\varrho)=-\varrho+cU_{\alpha_*}'(\varrho)=-\varrho+O(\varrho^{-1})$, contradicting
$\sigma<1$. This proves \eqref{eq:beff-O-s}.

Combining \eqref{eq:beff-O-s} with
\eqref{eq:Wtilde-refined}, we obtain
\begin{equation}
  |\widetilde W(z)|
  \le
  C\eta_\eps^2(1+|z|)^\sigma,
  \qquad
  |z|\le\frac{3r}{4s_\eps}.
  \label{eq:Wtilde-O2-weighted}
\end{equation}
Moreover, Lemma~\ref{lem:weighted-green}, together with the standard
interior gradient estimate for the harmonic boundary part, gives
\begin{equation}
  |\nabla\widetilde W(z)|
  \le
  C\eta_\eps^2,
  \qquad
  |z|\le\frac{3r}{4s_\eps}.
  \label{eq:grad-Wtilde-O2}
\end{equation}

\medskip
\noindent
{ 
\emph{Step 3: sharp first Fourier estimate.}
}
For $\Pi_1f(\varrho)=(2\pi)^{-1}\int_0^{2\pi}
f(\varrho \boldsymbol\omega(\theta))\boldsymbol\omega(\theta)\,d\theta$, set
\begin{equation}\label{eq:P1-definition}
 Z_{\rm eff}(z)=-s_\eps \mathbf{b}^{\rm eff}\cdot z
 -\frac{2s_\eps}{\kappa}\mathbf{b}^{\rm eff}\cdot\nabla U_\alpha(z),
 \qquad \mathcal R=\widetilde W-Z_{\rm eff}.
\end{equation}
Then $L_\alpha\mathcal R=E_\eps^{(2)}$ and
$\mathcal R(0)=\nabla\mathcal R(0)=0$.
To make the first-mode cancellation explicit, put
\[
 \begin{aligned}
 a^{(1)}(z)&=\Xi_1(z)+s_\eps\mathbf b\cdot z,
 &\Gamma^{(1)}(z)&=s_\eps\mathbf g\cdot z,\\
 \ell^{(2)}(z)&=\frac{s_\eps^2}{2}z^T\mathsf Bz,
 &\Gamma^{(2)}(z)&=\frac{s_\eps^2}{2}z^T\mathsf Cz.
 \end{aligned}
\]
where $\mathsf B,\mathsf C$ are defined in \eqref{eq:B-tracefree}.
The quadratic coefficient terms and the quadratic products of first
modes contribute $-\mathcal Q_2$ to $E_\eps^{(2)}$, where
\[
 \begin{aligned}
 \mathcal Q_2={}&F'(U)\ell^{(2)}+F(U)\Gamma^{(2)}
  +\frac12F''(U)(a^{(1)})^2 +F'(U)\Gamma^{(1)}a^{(1)}
  +\frac12F(U)(\Gamma^{(1)})^2.
 \end{aligned}
\]
At each fixed radius, $a^{(1)}$ and $\Gamma^{(1)}$ are pure first
Fourier modes, whereas $\ell^{(2)}$ and $\Gamma^{(2)}$ contain only
modes $0$ and $2$. Products of two first modes also contain only
modes $0$ and $2$, so $\Pi_1\mathcal Q_2=0$ exactly.
On $|z|\le3r/(4s_\eps)$, the remaining terms are estimated using
\[
 \begin{aligned}
 |a^{(1)}|+|\Gamma^{(1)}|\le Cs_\eps(1+|z|),\ \ \
 |\ell^{(2)}|+|\Gamma^{(2)}|\le Cs_\eps^2(1+|z|)^2,\ \ \
 |\widetilde W|\le Cs_\eps^2(1+|z|)^\sigma.
 \end{aligned}
\]
The coefficient Taylor remainders and cubic products give the power
$(1+|z|)^{3-m_*}$ below. The remaining mixed products with
$\widetilde W$ and its square give the powers
$(1+|z|)^{2+\sigma-m_*}$, $(1+|z|)^{1+\sigma-m_*}$, and
$(1+|z|)^{2\sigma-m_*}$; higher products are absorbed using
$s_\eps(1+|z|)\le C$ and $\sigma<1$.
Together with \eqref{eq:E2-crude} on the outer annulus, this yields
\begin{align}
 |\Pi_1E_\eps^{(2)}(\varrho)|
 &\le Cs_\eps^3\big[(1+\varrho)^{3-m_*}+(1+\varrho)^{2+\sigma-m_*}
       \notag\\[-1mm]
 &\hspace{32mm} +(1+\varrho)^{1+\sigma-m_*}+(1+\varrho)^{2\sigma-m_*}\big],
 &&\varrho\le\frac{3r}{4s_\eps},\label{eq:P1-error-after-refinement}\\
 |\Pi_1E_\eps^{(2)}(\varrho)|
 &\le Cs_\eps^2(1+\varrho)^{2-m_*},
 &&\frac{3r}{4s_\eps}\le \varrho\le\frac{r_\eps}{s_\eps},
 \label{eq:P1-error-outer}
\end{align}
where $r_\eps=r-|\mathbf{d}|$. Let $h_1=U_\alpha'$ and choose $h_2$ so that
$\mathcal W(h_1,h_2)=1/\varrho$; uniformly on $I$,
$|h_1|\le C/(1+\varrho)$ and $|h_2|\le C(\varrho+\varrho^{-1})$.
Variation of parameters gives
\begin{equation}\label{eq:P1-R-variation}
 \Pi_1\mathcal R(\varrho)
 =-h_1(\varrho)\int_0^\varrho h_2(\tau)\Pi_1E_\eps^{(2)}(\tau)\tau\,d\tau
 +h_2(\varrho)\int_0^\varrho h_1(\tau)\Pi_1E_\eps^{(2)}(\tau)\tau\,d\tau.
\end{equation}
Since $m_*>4$ and $\sigma<m_*-4$, these integrals satisfy
\begin{align}
 \int_0^{r_\eps/s_\eps}|h_1(\tau)|
       |\Pi_1E_\eps^{(2)}(\tau)|\tau\,d\tau
 &\le Cs_\eps^3,\label{eq:P1-growing-moment}\\
 |h_1(\varrho)|\int_0^\varrho|h_2(\tau)|
       |\Pi_1E_\eps^{(2)}(\tau)|\tau\,d\tau
 &\le Cs_\eps^2
 \quad(1\le \varrho\le r_\eps/s_\eps).
 \label{eq:P1-decaying-part}
\end{align}
The outer contribution in the first estimate is
$O(s_\eps^{m_*-1})=o(s_\eps^3)$. The growing term is also
$O(s_\eps^2)$ because $h_2(\varrho)=O(\varrho)$, and hence
\begin{equation}\label{eq:first-mode-remainder}
 |\Pi_1\mathcal R(\varrho)|\le Cs_\eps^2
 \qquad(1\le \varrho\le r_\eps/s_\eps).
\end{equation}
At $\varrho=r_\eps/s_\eps$, use $\mathbf{b}^{\rm eff}=O(s_\eps)$,
$U_\alpha'(\varrho)=O(s_\eps)$, and
$\Pi_1\Xi_1(\varrho)=O(s_\eps^2\de^2)$ to obtain
\begin{align}
 \Pi_1\widehat v(r_\eps/s_\eps)
 &=-\frac{r_\eps}{2}\mathbf{b}^{\rm eff}+O(s_\eps^2),
 \qquad 
 \Pi_1\widehat v(r_\eps/s_\eps)
 =-\frac m{2r}\mathbf{d}+O(\etae^2)=O(\etae^2).
 \label{eq:P1-green-boundary}
\end{align}
The second line is Lemma~\ref{lem:first-moment}. Comparing the two,
with $r_\eps=r+O(\etae^2)$ and $s_\eps\asymp\etae$, proves
$|\mathbf{b}^{\rm eff}|\le C\etae^2$, namely assertion~(ii).

\medskip\noindent
\emph{Step 4: the center shift.}
By \eqref{eq:Wtilde-O2-weighted}--\eqref{eq:grad-Wtilde-O2},
$\|W-\Xi_1\|_{C^1(B_R)}\le C_R\etae^2$ for every fixed $R$.
At $\zeta_\eps=-\mathbf d/s_\eps=O(\etae)$, the expansion
$Q_\alpha(\varrho)=O(\varrho^3)$ gives
$|\nabla\Xi_1(\zeta_\eps)|\le C\etae^3\de^2$, while
$\nabla U_\alpha(\zeta_\eps)=-\kappa\zeta_\eps/2+O(|\zeta_\eps|^3)$.
Insert these estimates and
$\nabla\mathfrak h(q)=\nabla\mathfrak h(q^*)+O(\etae^2)$
into \eqref{eq:max-condition-scaled} and multiply by $s_\eps\asymp\etae$:
\(
 \mathbf d=-\frac{2s_\eps^2}{\kappa}\nabla\mathfrak h(q^*)
             +O(\etae^3+\etae^4\de^2).
\)
This proves assertion~\textup{(iii)}.

\medskip
\noindent
\emph{Step 5: the second Fourier mode.}
For $\mathsf B,\mathsf C$ from \eqref{eq:B-tracefree}, Taylor expansion yields
\begin{align}
 \ell(z)=\log A(z)
 &=s_\eps \mathbf{b}\cdot z+\frac{s_\eps^2}{2}z^T\mathsf B z+O(s_\eps^3|z|^3),
 \label{eq:logA-second-correct}\\
 \Gamma(z)
 &=s_\eps \mathbf{g}\cdot z+\frac{s_\eps^2}{2}z^T\mathsf C z
   +O(s_\eps^3\de^3|z|^3).
 \label{eq:logGamma-second-correct}
\end{align}
The decomposition $z^T\mathsf S z=\tfrac12(\tr\mathsf S)|z|^2+z^T\mathsf S^\circ z$
separates modes $0$ and $2$, giving
\begin{align}
 e^\Gamma F(U+W+\ell)
 ={}&F(U)+F'(U)W+s_\eps[F'(U)\mathbf{b}+F(U)\mathbf{g}]\cdot z
 \notag\\
 &+\frac{s_\eps^2}{4}[F'(U)\tr\mathsf B+F(U)\tr\mathsf C]|z|^2
 \notag\\
 &+\frac{s_\eps^2}{2}[F'(U)z^T\mathsf B^\circ z+F(U)z^T\mathsf C^\circ z]
   +\mathcal R_\eps^{(2)}+\mathcal R_\eps^{(3)},
 \label{eq:full-second-order-expansion}
\end{align}
where $\mathbf{b}=O(\etae^2+\de^2)$, $W=\widetilde W+\Xi_1$, and the
established weighted bounds give
\begin{align}
 |\mathcal R_\eps^{(2)}|
 &\le C\big[\etae^4(1+|z|)^{-m_*+2\sigma}
       +\etae^4(1+|z|)^{-m_*+4}
 +(\etae^6+\etae^2\de^4)(1+|z|)^{-m_*+2}\big],
 \label{R1}\\
 |\mathcal R_\eps^{(3)}|
 &\le C\etae^3(1+|z|)^{-m_*+3}.
 \label{R2}
\end{align}
For either angular component of \eqref{eq:Psi2-equation}, let $f(\varrho)$
be its radial right-hand side. Let $h_0\sim\varrho^2$ at zero and
$h_\infty\sim\varrho^{-2}$ at infinity solve the homogeneous $k=2$ equation.
Their constant Wronskian factor
$\omega=\varrho(h_0h_\infty'-h_0'h_\infty)$ is nonzero by
Lemma~\ref{thm:linearized-classification}. The regular solution with
no growing branch is
\[
 \psi(\varrho)=\frac1\omega\left[
 h_\infty(\varrho)\int_0^\varrho t h_0(t)f(t)\,dt
 +h_0(\varrho)\int_\varrho^\infty t h_\infty(t)f(t)\,dt\right].
\]
Indeed, $f(\varrho)=O(\varrho^2)$ at zero and
$f(\varrho)=O(\varrho^{2-m_*})$ at infinity;
the formula gives $\psi=O(\varrho^2)$ at zero and bounded
$|\psi|+\varrho|\psi'|$ at infinity since $m_*>4$.
The bounds are uniform for $\alpha\in I$ by compactness and
$\omega\ne0$; the same kernel classification gives uniqueness.
Combining the two angular components defines $\Psi_2$ and gives
\begin{equation}\label{eq:Psi2-bound}
 \Psi_2(0)=0,\quad\nabla\Psi_2(0)=0,\qquad
 |\Psi_2(z)|+(1+|z|)|\nabla\Psi_2(z)|\le C.
\end{equation}
For \(
 \mathcal T=W-\Xi_1-s_\eps^2\Psi_2,\)
we obtain
\begin{align}
 L_\alpha\mathcal T
 ={}&-\frac{s_\eps^2}{4}[F'(U)\tr\mathsf B+F(U)\tr\mathsf C]|z|^2
 \notag\\
 &-s_\eps(\mathbf{b}^{\rm eff}\cdot z)F'(U)
   -\mathcal R_\eps^{(2)}-\mathcal R_\eps^{(3)},
 \qquad\mathcal T(0)=\nabla\mathcal T(0)=0.
 \label{eq:W-second-equation-correct}
\end{align}
Since $|\mathbf{b}^{\rm eff}|=O(\etae^2)$,
$|\tr\mathsf B|+|\tr\mathsf C|=O(\de^2)$, and $|z|\le C/\etae$,
\begin{equation}\label{eq:second-mode-equation-bound}
 |L_\alpha\mathcal T|
 \le C[\etae^3(1+|z|)^{-m_*+3}
       +\etae^2\de^2(1+|z|)^{-m_*+2}].
\end{equation}
Set $\tau_\eps=\etae^{2+\sigma}+\etae^2\de^2$ and choose
$0<\rho<m_*-4-\sigma$. The normalization of the first source term in
\eqref{eq:second-mode-equation-bound} is justified, uniformly for
$|z|\le C/\etae$, by
\[
 \etae^{1-\sigma}(1+|z|)^{5+\rho-m_*}
 \le C\bigl(\etae^{1-\sigma}
             +\etae^{m_*-4-\sigma-\rho}\bigr)\le C.
\]
Indeed, use $1+|z|\ge1$ when $5+\rho-m_*\le0$ and
$1+|z|\le C/\etae$ otherwise. Both powers of $\etae$ are positive,
since $\sigma<1$ and $m_*-4-\sigma-\rho>0$.
For the second source term,
\(\frac{\etae^2\de^2}{\tau_\eps} (1+|z|)^{4+\rho-m_*}\le1,\)
because $\tau_\eps\ge\etae^2\de^2$ and $4+\rho-m_*<0$.
Consequently,
$\tau_\eps^{-1}|L_\alpha\mathcal T|
\le C(1+|z|)^{-2-\rho}$, without any relation between $\etae$ and $\de$.
The harmonic extension $h_\eps$ of $\mathcal T|_{\partial D_\eps}$
has oscillation $O(\etae^2)$. Recall that
$D_\eps=B_{R_\eps}(\zeta_\eps)$, with
$R_\eps\asymp\etae^{-1}$ and $\zeta_\eps=O(\etae)$.
On $|z|\le R_\eps/2$, the interior gradient estimate gives
$|h_\eps(z)-h_\eps(0)|\le C\etae^3|z|$.
On the remaining part of $D_\eps$, the same inequality follows from
the oscillation bound and $|z|\ge R_\eps/2$.
Thus, using $\sigma<1$ and $\tau_\eps\ge\etae^{2+\sigma}$,
\[
 |h_\eps(z)-h_\eps(0)|\le C\etae^3|z|
 \le C\tau_\eps(1+|z|)^\sigma\qquad(z\in D_\eps).
\]
The full-disk Green and normalized-kernel argument of
Lemma~\ref{lem:preliminary-pointwise} now gives
$|\mathcal T|\le C\tau_\eps(1+|z|)^\sigma$.
This proves assertion~(iv).

\medskip\noindent
\emph{Step 6: the local mass and uniform estimate.}
By \eqref{eq:correct-source-change-variable},
$2\pi\widetilde m_{t,\eps}=\int_{D_\eps}\mathcal F_\eps$.
In \eqref{eq:full-second-order-expansion}, the first and trace-free
second modes integrate to zero on centered disks. The displacement
$\zeta_\eps=O(\etae)$ of $D_\eps$ contributes only a decaying tail;
indeed, the omitted radial mass is $O(\etae^{m_*-2})$.
The trace terms are $O(\etae^2\de^2)$, and the largest remainder
integrals are
\[
 \etae^3\int_0^{C/\etae}(1+\varrho)^{3-m_*}\varrho\,d\varrho
 +\etae^4\int_0^{C/\etae}(1+\varrho)^{4-m_*}\varrho\,d\varrho
 =O(\etae^{2+\sigma}).
\]
At $m_*=5$ or $6$ the additional logarithm is absorbed by the strict
choice \eqref{eq:sigma-choice}. Moreover,
$W=\mathcal T+\Xi_1+s_\eps^2\Psi_2$, so the same angular cancellation
and the weighted bound on $\mathcal T$ give
\[
 2\pi\widetilde m_{t,\eps}
 =2\pi m_{t,\eps}+\int_{D_\eps}F'(U_\alpha)W\,dz+O(\tau_\eps)
 =2\pi m_{t,\eps}+O(\tau_\eps).
\]
This proves assertion~(v).
Finally, on $|z|\le r/(2s_\eps)$, since $\sigma<1$,
\[
 \tau_\eps(1+|z|)^\sigma
 \le C(\etae^2+\etae^{2-\sigma}\de^2)
 \le C(\etae^2+\etae\de^2).
\]
Hence $|W|\le|\mathcal T|+|\Xi_1|+s_\eps^2|\Psi_2|
\le C(\etae^2+\etae\de^2)$, and interior elliptic estimates give
the same bound for $\nabla W$.
If $\varphi_i$ is constant, $\Xi_1$ and the metric trace vanish;
$\log H$ is already harmonic by its definition. This proves assertion~(i) and
completes the theorem.
\end{proof}

\subsection{Moment-sampling consequences}
\label{app:moment-consequences}

\begin{proposition}[Second-moment sampling formula]
\label{prop:moment-sampling}
Use the first-scale source $f_\eps$ from
\eqref{eq:first-scale-equation-sec3}. For
$f\in C^{2,\sigma}(\overline{B_r(q_{t,\eps})})$,
\begin{align}
  \int_{B_r(q_{t,\eps})}f(y)f_\eps(y)\,dy
  ={}&2\pi m_{t,\eps}f(q_{t,\eps}^*)
  -\etae^2J(\alpha_{t,\eps})
   \Delta f(q_{t,\eps}^*)
  \notag\\
  &+O\bigl(\etae^{2+\sigma}
  +\etae^2\de^2\bigr)
  \|f\|_{C^{2,\sigma}(B_r(q_{t,\eps}))}.
  \label{eq:moment-sampling}
\end{align}
In particular, the $\etae^2$ term vanishes whenever $f$ is harmonic in the
local disk.
\end{proposition}

\begin{proof}
Use \eqref{eq:correct-source-change-variable} on $D_{t,\eps}$.
Its zeroth moment is \eqref{eq:local-mass-refined}. The first two
weighted moments satisfy
\begin{align}
 \int_{D_{t,\eps}} z\,\mathcal F_{t,\eps}(z)\,dz
 &=O(\etae^{1+\sigma}+\etae\de^2),
 \label{eq:source-first-moment-correct}\\
 \int_{D_{t,\eps}} z_jz_k\,\mathcal F_{t,\eps}(z)\,dz
 &=\frac{\delta_{jk}}2\int_{\R^2}|z|^2F(U_{t,\eps})\,dz
   +O(\etae^\sigma(1+\de^2)),
 \label{eq:source-second-moment-correct}
\end{align}
Here is the tail and angular bookkeeping. For $j=0,1,2$, we have 
\(
 \int_{|z|\ge c/\etae}|z|^jF(U_{t,\eps})\,dz
 \le C\etae^{m_*-2-j}.
\)
Thus replacing the shifted disk by a centered one has admissible
error in each moment. In the first moment, the radial and second-mode
terms in \eqref{eq:full-second-order-expansion} vanish by parity.
The linear coefficient and $\Xi_1$ terms contribute
$O(\etae^3+\etae\de^2)$; the $\mathcal T$ term and
\eqref{R1}--\eqref{R2}, integrated with weight $|z|$, contribute
$O(\etae^{1+\sigma}+\etae\de^2)$.
For example the cubic term is bounded by
$C\etae^3\int_0^{C/\etae}(1+\varrho)^{4-m_*}\varrho\,d\varrho
=O(\etae^{1+\sigma})$.
With weight $|z|^2$, all perturbation terms are
$O(\etae^\sigma(1+\de^2))$; the slowest tail is
$\etae^{m_*-4}=O(\etae^\sigma)$.
These estimates prove \eqref{eq:source-first-moment-correct}--
\eqref{eq:source-second-moment-correct}; borderline logarithms are
again absorbed by \eqref{eq:sigma-choice}.

Taylor-expand $f(q_{t,\eps}^*+s_{t,\eps}z)$ to order two.
Its linear term is within the stated error by
\eqref{eq:source-first-moment-correct}. Radiality and
Lemma~\ref{lem:J-negative} together with  $s_{t,\eps}^2=\etae^2(1+O(\de^2))$ give the quadratic term
\[
 \frac{s_{t,\eps}^2}{4}\Delta f(q_{t,\eps}^*)
       \int_{\R^2}|z|^2F(U_{t,\eps})\,dz
 =-\etae^2J(\alpha_{t,\eps})\Delta f(q_{t,\eps}^*)
   +O(\etae^2\de^2)\|f\|_{C^{2,\sigma}}.
\]
The Taylor remainder is integrable by
$2+\sigma<m_*-2$. This proves \eqref{eq:moment-sampling}.
\end{proof}


\section*{Acknowledgments}

The authors are grateful to Professor Ky Ho of the University of Economics
Ho Chi Minh City for valuable discussions throughout the development
and completion of this work.

The authors also acknowledge the use of OpenAI's Codex AI assistant
in developing the construction and its proof, checking algebraic
identities, and organizing and preparing the \LaTeX{} manuscript.
The authors are responsible for the mathematical arguments and the
final content of the manuscript.


\section*{Statements and Declarations}



\subsection*{Funding}
Y. Lee was supported by a National Research Foundation of Korea (NRF) grant funded by the Korea government (MSIT) (No. RS-2026-25489601). L. Zhang was supported by a Simons foundation travel grant Award ID:SFI-MPS-TSM-00013752. 

\subsection*{Competing interests}
The authors declare that they have no competing interests.

\subsection*{Data availability}
No datasets were generated or analysed in this theoretical study.
Data sharing is therefore not applicable to this article.

\begin{thebibliography}{99}

\bibitem{BartolucciTarantello}
D.~Bartolucci and G.~Tarantello,
\emph{Asymptotic blow-up analysis for singular Liouville type equations
with applications},
J. Differential Equations \textbf{262} (2017), no.~7, 3887--3931.
\url{https://doi.org/10.1016/j.jde.2016.12.003}.

\bibitem{BM}
H.~Brezis and F.~Merle,
\emph{Uniform estimates and blow-up behavior for solutions of
$-\Delta u=V(x)e^u$ in two dimensions},
Comm. Partial Differential Equations \textbf{16} (1991), 1223--1253.
\url{https://doi.org/10.1080/03605309108820797}.

\bibitem{CaffarelliYang}
L.~A.~Caffarelli and Y.~Yang,
\emph{Vortex condensation in the Chern--Simons Higgs model:
An existence theorem},
Comm. Math. Phys. \textbf{168} (1995), no.~2, 321--336.
\url{https://doi.org/10.1007/BF02101552}.

\bibitem{CFL}
H.~Chan, C.-C.~Fu and C.-S.~Lin,
\emph{Non-topological multi-vortex solutions to the self-dual
Chern--Simons--Higgs equation},
Comm. Math. Phys. \textbf{231} (2002), 189--221.
\url{https://doi.org/10.1007/s00220-002-0691-6}.

\bibitem{Choe}
K.~Choe,
\emph{Asymptotic behavior of condensate solutions in the
Chern--Simons--Higgs theory},
J. Math. Phys. \textbf{48} (2007), 103501.

\bibitem{ChoeKim}
K.~Choe and N.~Kim,
\emph{Blow-up solutions of the self-dual Chern--Simons--Higgs vortex
equation},
Ann. Inst. H. Poincar\'e Anal. Non Lin\'eaire \textbf{25} (2008), 313--338.
\url{https://doi.org/10.1016/j.anihpc.2006.11.012}.

\bibitem{CKL}
K.~Choe, N.~Kim and C.-S.~Lin,
\emph{Existence of self-dual non-topological solutions in the
Chern--Simons Higgs model},
Ann. Inst. H. Poincar\'e Anal. Non Lin\'eaire \textbf{28} (2011), 837--852.
\url{https://doi.org/10.1016/j.anihpc.2011.06.003}.

\bibitem{DAprileWeiZhangAJM}
T.~D'Aprile, J.~Wei and L.~Zhang,
\emph{On non-simple blowup solutions of Liouville equation},
Amer. J. Math., to appear.
\url{https://arxiv.org/abs/2209.05271}.

\bibitem{DAprileWeiZhangCVPDE}
T.~D'Aprile, J.~Wei and L.~Zhang,
\emph{On the construction of non-simple blow-up solutions for the singular
Liouville equation with a potential},
Calc. Var. Partial Differential Equations \textbf{63} (2024), Paper No.~68.
\url{https://doi.org/10.1007/s00526-024-02676-x}.

\bibitem{DAprileWeiZhangDCDS}
T.~D'Aprile, J.~Wei and L.~Zhang,
\emph{Non-simple blow-up for singular Liouville equations in unit ball},
Discrete Contin. Dyn. Syst. \textbf{44} (2024), no.~7, 1938--1957.
\url{https://doi.org/10.3934/dcds.2024015}.

\bibitem{DAprileWeiZhangTAMS}
T.~D'Aprile, J.~Wei and L.~Zhang,
\emph{High order vanishing theorems for non-simple blowup solutions of
singular Liouville equation},
Trans. Amer. Math. Soc. (2026), published electronically August 20, 2026.
\url{https://doi.org/10.1090/tran/9742}.

\bibitem{DelPinoEspositoFigueroaMusso}
M.~del Pino, P.~Esposito, P.~Figueroa, and M.~Musso,
\emph{Nontopological condensates for the self-dual Chern--Simons--Higgs model},
Comm. Pure Appl. Math. \textbf{68} (2015), no.~7, 1191--1283.
\url{https://doi.org/10.1002/cpa.21548}.

\bibitem{HongKimPac}
J.~Hong, Y.~Kim and P.-Y.~Pac,
\emph{Multivortex solutions of the Abelian Chern--Simons--Higgs theory},
Phys. Rev. Lett. \textbf{64} (1990), 2230--2233.
\url{https://doi.org/10.1103/PhysRevLett.64.2230}.

\bibitem{JackiwWeinberg}
R.~Jackiw and E.~J.~Weinberg,
\emph{Self-dual Chern--Simons vortices},
Phys. Rev. Lett. \textbf{64} (1990), 2234--2237.
\url{https://doi.org/10.1103/PhysRevLett.64.2234}.

\bibitem{KuoLin}
T.-J.~Kuo and C.-S.~Lin,
\emph{Estimates of the mean field equations with integer singular sources:
non-simple blowup},
J. Differential Geom. \textbf{103} (2016), 377--424.
\url{https://doi.org/10.4310/jdg/1468517500}.

\bibitem{LiLiu}
J.~Y.~Li and L.~Liu,
\emph{The qualitative behavior at a vortex point for the Chern--Simon--Higgs equation},
arXiv:2306.03687 (2023).
\url{https://arxiv.org/abs/2306.03687}.

\bibitem{LinYanCMP}
C.-S.~Lin and S.~Yan,
\emph{Bubbling solutions for relativistic Abelian Chern--Simons model on a torus},
Comm. Math. Phys. \textbf{297} (2010), 733--758.
\url{https://doi.org/10.1007/s00220-010-1056-1}.

\bibitem{LinYanARMA}
C.-S.~Lin and S.~Yan,
\emph{Existence of bubbling solutions for Chern--Simons model on a torus},
Arch. Ration. Mech. Anal. \textbf{207} (2013), 353--392.
\url{https://doi.org/10.1007/s00205-012-0575-7}.

\bibitem{LinYanAIHP}
C.-S.~Lin and S.~Yan,
\emph{On condensate of solutions for the Chern--Simons--Higgs equation},
Ann. Inst. H. Poincar\'e Anal. Non Lin\'eaire \textbf{34} (2017), 1329--1354.
\url{https://doi.org/10.1016/j.anihpc.2016.10.006}.

\bibitem{LinYanAdvMath}
C.-S.~Lin and S.~Yan,
\emph{On the mean field type bubbling solutions for Chern--Simons--Higgs equation},
Adv. Math. \textbf{338} (2018), 1141--1188.
\url{https://doi.org/10.1016/j.aim.2018.09.021}.

\bibitem{FredholmNotes}
S.~Nonnenmacher and M.~L\'eautaud,
\emph{Introduction to Spectral Theory},
lecture notes, Universit\'e Paris-Saclay, October 17, 2025,
Theorem~5.3.1.
\url{https://leautaud.perso.math.cnrs.fr/files/spectral-theory-2025.pdf}.

\bibitem{SpruckYang}
J.~Spruck and Y.~Yang,
\emph{The existence of non-topological solitons in the self-dual
Chern--Simons theory},
Comm. Math. Phys. \textbf{149} (1992), 361--376.
\url{https://doi.org/10.1007/BF02097630}.

\bibitem{Tarantello}
G.~Tarantello,
\emph{Multiple condensate solutions for the Chern--Simons--Higgs theory},
J. Math. Phys. \textbf{37} (1996), no.~8, 3769--3796.
\url{https://doi.org/10.1063/1.531601}.

\bibitem{WeiZhang}
J.~Wei and L.~Zhang,
\emph{Estimates for Liouville equation with quantized singularities},
Adv. Math. \textbf{380} (2021), Paper No.~107606.
\url{https://doi.org/10.1016/j.aim.2021.107606}.

\bibitem{WeiZhangVanishing}
J.~Wei and L.~Zhang,
\emph{Vanishing estimates for Liouville equation with quantized
singularities},
Proc. Lond. Math. Soc. (3) \textbf{124} (2022), no.~1, 106--131.
\url{https://doi.org/10.1112/plms.12426}.

\bibitem{WeiZhangLaplacian}
J.~Wei and L.~Zhang,
\emph{Laplacian vanishing theorem for a quantized singular Liouville
equation},
J. Eur. Math. Soc. \textbf{28} (2026), no.~1, 237--267.
\url{https://doi.org/10.4171/JEMS/1482}.

\end{thebibliography}
\end{document}
